% !TEX root = main.tex
% Appendix. A: executors. B: theorem arguments followed by technical verification.
% C: numerical verification of the theory. D: experiments by prediction. E: reproducibility.
\providecommand{\Zk}{\mathbb Z_K}
\providecommand{\reg}[1]{\ensuremath{\mathsf{#1}}}
\providecommand{\op}[1]{\textsc{#1}}

\section*{Roadmap}
\addcontentsline{toc}{section}{Roadmap}

The table maps results to assumptions, proofs and numerical checks.
\Cref{app:experiments} gives the training protocols and results.
\Cref{app:reproducibility} lists the code and compute.

\begin{center}\small
\fitwidth{\begin{tabular}{@{}llll@{}}
\toprule
Result & Assumptions & Proof & Numerical check\\
\midrule
\Cref{eq:grad-split}: credit matches degree & independent moves & \cref{app:tools} & exact gradients, \cref{app:exp-spectra}\\
\Cref{eq:horizon}: every unseen move costs $\mixrate$ & (A3) & \cref{lem:target-contraction} & \cref{tab:verify}, row 1\\
\Cref{prop:represent}: both executors in two blocks & (A1) & \cref{app:exact} & \cref{tab:verify}, row 8\\
\Cref{thm:main}: outcome interaction credit & (A1)--(A3), $0<\mixrate<1$ & \cref{app:thm1} & \cref{tab:verify}, rows 2--3\\
\Cref{prop:plateau}: plateau with a nonzero readout & (A1)--(A4), $0<\mixrate<1$ & \cref{app:thm1} & \cref{tab:verify}, row 4; \cref{app:exp-regime}\\
\Cref{thm:separation}: horizon-free gradient and rule credit & (A1)--(A3), uniform start & \cref{app:thm2} & \cref{tab:verify}, rows 5--7; \cref{app:exp-regime}\\
\Cref{thm:outcome-end}: outcome never reaches the basin & (A1)--(A3), $\initscale\le c\mixrate$, $\tilde V=0$ & \cref{app:end} & --\\
\Cref{thm:basin}, \cref{cor:exact}: convergence in the basin & Armijo steps, \cref{ass:armijo} & \cref{app:end} & --\\
\bottomrule
\end{tabular}}
\end{center}

%=============================================================================
\section{Exact two-block executors}
\label{app:architectures}

We prove \cref{prop:represent} and list the nonzero parameters.

\subsection{Program overview}

The table specifies the operations in each block.
$\widetilde s_{t-1}$ denotes the displayed previous state.
All state arithmetic is modulo $K$.

\begin{center}\small
\begin{tabular}{@{}llll@{}}
\toprule
Program part & Reads & Writes & Purpose\\
\midrule
Block 1, head & query token & $\reg{qry}$ & copy the query to letter positions\\
Block 1, MLP & letter, $\reg{qry}$ & $\reg{match}$ & compute $b_t$\\
\midrule
Step, block 2, head 1 & previous value & $\reg{prev}$ & copy $\widetilde s_{t-1}$\\
Step, block 2, head 2 & $\reg{match}$ at letter $t$ & $\reg{mt}$ & copy $b_t$\\
Step, block 2, MLP & $\reg{prev}$, $\reg{mt}$ & $\reg{out}$ & compute $s_t$\\
\midrule
Outcome, block 2, head 1 & $\reg{match}$ at all letters & $\reg{cnt}$ & compute $m=\sum_tb_t$\\
Outcome, block 2, head 2 & start value & $\reg{s0}$ & copy $s_0$\\
Outcome, block 2, MLP & $\reg{s0}$, $\reg{cnt}$ & $\reg{out}$ & compute $s_n$\\
\bottomrule
\end{tabular}
\end{center}

A \emph{route} copies a register from one source position.
An \emph{average} reads several sources.
A \emph{lookup} maps an input pair to its output.

\subsection{Architecture, layout and primitives}
\label{app:construction}
\label{app:model}

\textbf{Architecture.} Each block has three heads and one MLP, without LayerNorm.
For this construction, the head dimension $d$ equals the number of positions:

\begin{align*}
 h_i&\leftarrow h_i+\sum_{h=1}^{3}\sum_{j\le i}
 \operatorname{softmax}_{j\le i}\!\Big(\tfrac{1}{\sqrt d}(W^{(h)}_Qh_i)^{\!\top}W^{(h)}_Kh_j\Big)\,W^{(h)}_Vh_j,\\
 h_i&\leftarrow h_i+W_{\mathrm{out}}\relu(W_{\mathrm{in}}h_i+b_{\mathrm{in}}),
 \qquad
 \mathrm{logits}_i=V\,h_i .
\end{align*}

\textbf{Registers.} The residual stream is split into the registers of the table, in this order,
with $d=2n+6$ positions and MLP width $\max\{|\Sigma|,3K(n+1),2n+1\}$:
\[
 \underbrace{\reg{tl}}_{|\Sigma|}\;\underbrace{\reg{tv}}_{K}\;\underbrace{\reg{pos}}_{d}\;
 \underbrace{\reg{qry}}_{|\Sigma|}\;\underbrace{\reg{match}}_{1}\;\underbrace{\reg{prev}}_{K}\;
 \underbrace{\reg{mt}}_{1}\;\underbrace{\reg{cnt}}_{1}\;\underbrace{\reg{s0}}_{K}\;
 \underbrace{\reg{out}}_{K},
 \qquad d_{\mathrm{res}}=2|\Sigma|+4K+d+3 .
\]
$\reg{tl}$ and $\reg{tv}$ hold the token (a letter or a value) and $\reg{pos}$ the position, all
one-hot. $\reg r[a]$ is coordinate $a$ of register $\reg r$, and $E_{\reg a\leftarrow\reg b}$ is the
$d_{\mathrm{res}}\times d_{\mathrm{res}}$ matrix that copies register $\reg b$ into register $\reg a$.

\textbf{Tokens and positions.} Letters are $0,\ldots,|\Sigma|-1$, value $v$ is $|\Sigma|+v$, and SEP,
``:'' and EOS follow. Positions are $0$-indexed:
\begin{align*}
 \text{prompt: }&\ s_0@0,\quad q@1,\quad w_t@(1+t),\quad \text{SEP}@(n+2),\\
 \text{step: }&\ s_t@(n+2+t),\quad \text{``:''}@(2n+3),\quad s_n@(2n+4),\\
 \text{outcome: }&\ \text{``:''}@(n+3),\quad s_n@(n+4).
\end{align*}
Position $i$ predicts the token at $i+1$, so $n+1+t$ predicts $s_t$ in the step row and $n+3$
predicts $s_n$ in the outcome row.

\textbf{Route and average.} Give source positions score $\beta_n$ and all others score $0$.
Let $\mathrm{src}(i)\subseteq\{0,\ldots,i\}$ be the source set.
Unlisted positions use the fallback $\{f\}$.
Set $C_Q=\beta_n\sqrt d$ and $\beta_n=\max\{16,3\log(8(n+3))\}$:

\begin{equation}
 \begin{aligned}
 (W_K)_{j,\reg{pos}[j]}=1,\quad
 (W_Q)_{j,\reg{pos}[i]}&=C_Q\quad(j\in\mathrm{src}(i)),\\
 \tfrac1{\sqrt d}(W_Qh_i)^\top W_Kh_j&=\beta_n\,[j\in\mathrm{src}(i)].
 \end{aligned}
 \label{eq:route}
\end{equation}
With $W_V=E_{\reg a\leftarrow\reg b}$, the head writes the average of register $\reg b$ over $\mathrm{src}(i)$
into $\reg a$, up to a small leak. One source is a route and several are an average. A fallback
$f>i$ gives uniform attention at a position that predicts nothing supervised.

\textbf{Lookup.} Use one unit per input pair.
For one-hot inputs, let $\mathrm{pre}_{s'x'}=[s=s']+[x=x']-3/2$.
Only the matching pair has positive pre-activation, equal to $1/2$.
Output weights $2\bm e_{F(s',x')}$ give

\begin{equation}
 \textstyle\sum_{s',x'}2\,\bm e_{F(s',x')}\relu(\mathrm{pre}_{s'x'})=\bm e_{F(s,x)}.
 \label{eq:lookup}
\end{equation}

\subsection{Both programs are exact}
\label{app:exact}


\begin{proof}[Proof of \cref{prop:represent}]
\proofstep{Ideal routing.}
The parameters in \cref{app:params} give
\begin{align*}
 \reg{match}_{1+t}
 &=2\sum_{a\in\Sigma}\relu([w_t=a]+[q=a]-3/2)=b_t,\\
 \reg{prev}_{n+1+t}&=\bm e_{s_{t-1}},\qquad
 \reg{mt}_{n+1+t}=b_t,\\
 2\relu([s_{t-1}=j]+[b_t=\iota]-3/2)
 &=[(s_{t-1},b_t)=(j,\iota)],\\
 \reg{out}_{n+1+t}
 &=\sum_{j,\iota}\bm e_{j+\iota}[(s_{t-1},b_t)=(j,\iota)]
 =\bm e_{s_t}.
\end{align*}
For the outcome program,
\begin{align*}
 \reg{cnt}_{n+3}&=n\left(\frac1n\sum_{t=1}^nb_t\right)=m,
 \qquad\reg{s0}_{n+3}=\bm e_{s_0},\\
\end{align*}
\[
 \relu(m-a+1)-2\relu(m-a)+\relu(m-a-1)=[m=a].
\]
\begin{align*}
 \reg{out}_{n+3}&=\sum_{j,a}\bm e_{j+a}[s_0=j][m=a]\\
 &=\bm e_{s_0+m}=\bm e_{s_n}.
\end{align*}
All state indices are modulo $K$; the readout in \cref{app:params} supplies formatting and the final answer.

\proofstep{Finite-score error.}
With $d=2n+6$ and $\beta_n=\max\{16,3\log(8(n+3))\}$,
\[
 \rho_n=(d-1)e^{-\beta_n}
 <\frac{2(n+3)}{(8(n+3))^3}
 =\frac1{256(n+3)^2}.
\]
The routing and lookup errors satisfy
\begin{align*}
 |\reg{match}-b_t|&\le2\rho_n,\\
 \|\reg{prev}-\bm e_{s_{t-1}}\|&\le\rho_n,
 \quad |\reg{mt}-b_t|\le3\rho_n,\\
 \mathrm{pre}_{\mathrm{correct}}&\ge\tfrac12-4\rho_n>0>\mathrm{pre}_{\mathrm{other}},
 \quad\reg{out}_{s_t}\ge1-8\rho_n,\\
 |\reg{cnt}-m|&\le3n\rho_n,
 \quad\|\reg{s0}-\bm e_{s_0}\|\le\rho_n,\\
 e_{\mathrm{leak}}&=3n\rho_n+(n+2)\rho_n<\tfrac1{192},\\
 \reg{out}_{s_n}&\ge1-e_{\mathrm{leak}},
 \quad\reg{out}_{\mathrm{other}}\le e_{\mathrm{leak}}.
\end{align*}
The three-ReLU equality test is $1$-Lipschitz. Thus every supervised position has margin
\[
 z_{\mathrm{correct}}-\max_{a\ne\mathrm{correct}}z_a
 \ge\min\{20(1-8\rho_n),20(1-2e_{\mathrm{leak}})\}>19.
\]

\proofstep{Open neighbourhood.}
For fixed $n$, let $\mathcal P_n$ be the finite set of supported gold prefixes and prediction positions.
Continuity and the preceding margin give $\ballradius_{\mathrm{exec},n}>0$ such that
\[
 \|\theta-\theta^\star\|<\ballradius_{\mathrm{exec},n}
 \quad\Longrightarrow\quad
 \max_{P\in\mathcal P_n}\|\bm z_\theta(P)-\bm z_{\theta^\star}(P)\|_\infty<19/2.
\]
Every greedy prediction remains correct. Induction over the continuation proves the claim.
\end{proof}


\subsection{Explicit parameters}
\label{app:params}

All unspecified entries are zero.
The reconstruction script checks these tensors against the implementation entry by entry
(\cref{app:reproducibility}).

\textbf{Embeddings.}
$(W_{\mathrm{te}})_{a,\reg{tl}[a]}=1$, $(W_{\mathrm{te}})_{|\Sigma|+v,\reg{tv}[v]}=1$ and $(W_{\mathrm{pe}})_{i,\reg{pos}[i]}=1$.

\textbf{Block 1 (shared): the match bit.} A route copies the query letter to each letter
position, and a lookup with one unit per letter tests whether the two letters agree:
\begin{align*}
 \text{head 1: }&\ \mathrm{src}(1+t)=\{1\}\ (t=1,\ldots,n),\quad f=0,\quad W_V=E_{\reg{qry}\leftarrow\reg{tl}},\\
 \text{MLP unit }a:&\ (W_{\mathrm{in}})_{a,\reg{tl}[a]}=(W_{\mathrm{in}})_{a,\reg{qry}[a]}=1,\quad
 (b_{\mathrm{in}})_a=-\tfrac32,\quad (W_{\mathrm{out}})_{\reg{match},a}=2,
\end{align*}
so $\reg{match}_{1+t}=2\sum_a\relu([w_t=a]+[q=a]-\frac32)=b_t$. Heads 2 and 3 are zero, and non-letter
positions get $\reg{qry}=\reg{match}=0$.

\textbf{Step executor, block 2.} Two routes bring the previous state and the current bit to the
position predicting $s_t$, and a third lets the answer position copy $s_n$:
\begin{align*}
 \text{head 1: }&\ \mathrm{src}(n+2)=\{0\},\ \mathrm{src}(n+1+t)=\{n+1+t\}\ (t=2,\ldots,n),\ f=1,\\
 &\ W_V=E_{\reg{prev}\leftarrow\reg{tv}},\\
 \text{head 2: }&\ \mathrm{src}(n+1+t)=\{1+t\}\ (t=1,\ldots,n),\ f=0,
 \quad W_V=E_{\reg{mt}\leftarrow\reg{match}},\\
 \text{head 3: }&\ \mathrm{src}(2n+3)=\{2n+2\},\ f=1,
 \quad W_V=E_{\reg{out}\leftarrow\reg{tv}}.
\end{align*}
The MLP is the lookup \cref{eq:lookup} for $F(s,b)=s+b$, with $2K$ units $u=2j+\iota$ for $j\in[K]$ and
$\iota\in\{0,1\}$:
\[
 (W_{\mathrm{in}})_{u,\reg{prev}[j]}=1,\quad
 (W_{\mathrm{in}})_{u,\reg{mt}}=2\iota-1,\quad
 (b_{\mathrm{in}})_u=-\tfrac12-\iota,\quad
 (W_{\mathrm{out}})_{\reg{out}[j+\iota\bmod K],\,u}=2 .
\]

\textbf{Outcome executor, block 2.} An average over the letter positions, scaled by $n$, gives
the count $m$, and a route copies $s_0$:
\begin{align*}
 \text{head 1: }&\ \mathrm{src}(n+3)=\{2,\ldots,n+1\},\ f=0,\quad W_V=n\,E_{\reg{cnt}\leftarrow\reg{match}},\\
 \text{head 2: }&\ \mathrm{src}(n+3)=\{0\},\ f=1,\quad W_V=E_{\reg{s0}\leftarrow\reg{tv}}.
\end{align*}
The count $m$ is a scalar. Three ReLUs test equality to $a$:
$[m=a]=\relu(m-a+1)-2\relu(m-a)+\relu(m-a-1)$.
A gate $b_{\mathrm{gate}}=n+2$ disables the test unless $s_0=j$.
Use $3K(n+1)$ units $u=3((n+1)j+a)+u_{\mathrm{hinge}}$, for $j\in[K]$, $a\in\{0,\ldots,n\}$ and
$(u_{\mathrm{hinge}},\chi,w_\chi)\in\{(0,1,1),(1,0,-2),(2,-1,1)\}$:

\begin{gather*}
 (W_{\mathrm{in}})_{u,\reg{cnt}}=1,\qquad
 (W_{\mathrm{in}})_{u,\reg{s0}[j]}=b_{\mathrm{gate}},\\
 (b_{\mathrm{in}})_u=\chi-a-b_{\mathrm{gate}},\\
 (W_{\mathrm{out}})_{\reg{out}[j+a\bmod K],\,u}=w_\chi .
\end{gather*}
The three units of $(j,a)$ add to $[s_0=j]\,[m=a]$, so $\reg{out}=\bm e_{s_0+m}=\bm e_{s_n}$. The table has one
triple per pair $(s_0,m)$, which is why it grows with $n$.

\textbf{Readout.} $V_{|\Sigma|+v,\reg{out}[v]}=20$ and $V_{\text{``:''},\reg{pos}[i_:]}=V_{\mathrm{EOS},\reg{pos}[i_E]}=20$, with
$(i_:,i_E)=(2n+2,2n+4)$ for the step program and $(n+2,n+4)$ for the outcome program.

%=============================================================================
\section{Proofs of the theorems}
\label{app:damping-proof}
\label{app:mainproofs}

\paragraph{Proof map.}
The outcome theorem pairs the target bound in \cref{lem:target-contraction} with the
Jacobian bound in \cref{prop:score-degree}, then sums the credit.
The step theorem reuses the model bound.
Its local signal instead requires the current move and a state, with a position-selection
cost of order $1/n$.
\Cref{app:end} proves the end-of-training results.
\Cref{app:network} verifies the Transformer bounds after these theorem arguments.

\subsection{Common assumptions and degree calculus}
\label{app:assumptions}

(A1)--(A4) restate \cref{ass:arch,ass:init,ass:task,ass:gd} with the details the proofs use.

\begin{enumerate}[label=(A\arabic*),leftmargin=*,itemsep=1pt,topsep=2pt]
\item \emph{Architecture.} Use the two-block causal Transformer of \cref{app:model}.
Each block has three softmax heads and a ReLU MLP.
Token and position embeddings are learned. There is no LayerNorm.
A linear readout produces $K$ value logits $\bm z_\theta$.
Widths are fixed; constants may depend on them.
The executors instead use widths growing with $n$.

\item \emph{Initialization.} Let $\theta_0=\initscale\tilde\theta$ with fixed $\tilde\theta$ and $\initscale\le\initscale_0$.
Every MLP bias of $\tilde\theta$ has absolute value at least $1$.
Its value-token embeddings are linearly independent.
Write $\varsigma$ for their smallest singular value.
The bounds below hold on every ball of radius $\ballradius'\initscale$ around $\theta_0$ with
$\ballradius'\le\min(1/4,\varsigma/2)$. We fix $\ballradius=\min(1/8,\varsigma/4)$ and
$\initball=\{\theta:\|\theta-\theta_0\|\le\ballradius\initscale\}$, so the ball of radius $2\ballradius\initscale$ is covered too.

\item \emph{Task.} A permutation task with $y=s_n$ and contraction $\mixrate$, with moves that are
independent and identically distributed given the query.
\item[(A3$'$)] \emph{Uniform start.} The start state is uniform.
\item \emph{Training.} Full-batch gradient descent $\theta_{j+1}=\theta_j-\stepsize\,\pd{L}{\theta}(\theta_j)$ with any step size $\stepsize$;
$j\stepsize$ is the training time after $j$ updates. $\Lout$ is the cross-entropy of the
answer. For teacher-forced state-token losses $\ell_t$,
\[
 \Lstep=\frac1n\sum_{t=1}^n\E\ell_t,\qquad
 \Lstep^{\ge2}=\frac1{n-1}\sum_{t=2}^n\E\ell_t.
\]
The step-credit theorem and its quantitative checks use $\Lstep^{\ge2}$.
Formatting tokens and the final answer position are excluded from both state losses.
\item \emph{Programs.} $\thstep$ and $\thout$ are the exact step and outcome executors of
\cref{prop:represent}.
\end{enumerate}

\paragraph{Contraction.} Write $\bm T_x$ for the permutation matrix of move $x$, so that
$\bm e_{s_t}^\top=\bm e_{s_{t-1}}^\top\bm T_{x_t}$. With the query fixed,
\begin{equation}
 \overline{\bm T}:=\mathbb E_x\bm T_x,
 \qquad
 \mixrate:=\max_{\bm v\perp\one,\ \|\bm v\|=1}\|\bm v^\top\overline{\bm T}\|,
 \label{eq:lambda}
\end{equation}
maximized over queries if the move law depends on the query. Since $\|\bm v^\top\overline{\bm T}\|\le\mathbb E_x\|\bm v^\top\bm T_x\|=\|\bm v\|$,
always $\mixrate\le1$. With finitely many moves, $\mixrate<1$ if and only if, for every query, the products
$\bm T_x\bm T_{x'}^{-1}$ generate a group acting transitively on the states \citep{levin2017markov}.



\begin{remark}[The bias condition]
\label{rem:two-block-bias}
The bias condition in (A2) fixes one ReLU activation pattern for all inputs on $\initball$. The AdamW
runs violate it: at their initialization ($n=8$, MLP width $54$, $\initscale=0.02$, Gaussian biases), $10$ to
$19$ of the $54$ units per block switch with the letters. This is why those runs lie outside the
regime of the theorems, and why \cref{app:exp-regime} repeats the key measurements inside it.
\end{remark}

% Generated by tools/generate_notation_table.py. Edit notation.tex.
\begin{center}\small
\fitwidth{\begin{tabular}{@{}ll@{}}
\toprule
Symbol & Meaning\\
\midrule
$n,K,\Sigma$ & horizon, number of states, alphabet\\
$t,j,i,r,k$ & sequence index, update, item/position, degree, score count\\
$\theta,\tilde\theta,\initscale$ & parameters, initialization direction, initialization scale; \cref{ass:init}\\
$\ballradius,\initball$ & relative ball radius and initialization ball; \cref{ass:init}\\
$\failprob,\stepsize$ & failure probability and step size\\
$s_t,x_t,b_t,q$ & state, move, match bit, query; \cref{sec:setting}\\
$p,\omega$ & match probability and prompt/context\\
$\mixrate$ & contraction per unseen move; \cref{def:lambda}\\
$S,(f)_r$ & move-support set and degree-$r$ component; \cref{def:degree}\\
$\Lout,\Lstep,\Lstep^{\ge2}$ & outcome and state losses; \cref{def:losses}\\
$\Lstep^{\mathrm{full}}$ & loss on every generated step-format token; \cref{cor:exact}\\
$\cred_g,\outcredit,\rulecredit$ & credit functionals; \cref{def:credit,sec:steps}\\
$\rulecred,g_t$ & rule-credit norm and rule signal; \cref{sec:steps}\\
$A_r,J_r$ & target and Jacobian spectra; \cref{eq:spectra}\\
$\Jtotal,\Jhigh$ & total energy $\Jtotal=\sum_r J_r$ and higher-degree energy $\Jhigh=\sum_{r\ge2}J_r$\\
$J^{[k]}$ & Jacobian coefficient with $k$ scores; \cref{prop:score-degree}\\
$V,h,\bm z,\bm u$ & readout, feature, logits, uniform token distribution\\
$\items,\pi_i,\pi_\omega,\pimin$ & supervised items, item/prompt weights, smallest item weight; \cref{def:basin}\\
$\basin(L),\entryidx$ & loss basin and basin-entry update; \cref{def:basin,thm:basin}\\
$\margin_i,\margin_{\entryidx},\execmargin$ & item, entry-time, and executor logit margins\\
$\traintime_j,\traintime(\ballradius)$ & accumulated training time and ball-exit time\\
$\thstep,\thout$ & exact counting executors; \cref{prop:represent}\\
$\bm T_x,\overline{\bm T}$ & move permutation and its mean; \cref{eq:lambda}\\
$\xi,\zeta,\mu_j$ & score scale, Fourier root of unity, Fourier eigenvalue\\
$m_i,l_t,\kappa_{l,i}$ & visible-position count, current-letter position, positional dilution\\
$D,M_\ell,\bm o_\ell$ & fixed ReLU pattern, affine block map, bias offset; \cref{app:network}\\
$\partial_v,\partial^k$ & directional and higher differentials; \cref{app:network}\\
$\attconst,c,C$ & constants depending only on widths, $K$, and $\tilde\theta$\\
$\Ktokens,\blocks$ & number of output tokens and Transformer blocks\\
\bottomrule
\end{tabular}}
\captionof{table}{\textbf{Notation.} Symbols used in one proof only are defined there.
Named subscripts are roman; indices are italic.}
\label{tab:notation}
\end{center}


\paragraph{Degree calculus.}
\label{app:tools}

\paragraph{Efron--Stein decomposition \citep{efron1981jackknife,odonnell2014analysis}.}
For independent $x_1,\ldots,x_n$, a square-integrable $f$ has the unique decomposition
$f=\sum_S f_S$, where
$f_S=\sum_{U\subseteq S}(-1)^{|S\setminus U|}\E[f\mid x_U]$.
Each component depends only on $x_S$ and is centred in each move in $S$.
Orthogonality gives $\E[f_S^\top g_{S'}]=0$ for $S\ne S'$.
If $f$ depends only on $x_U$, then $f_S=0$ unless $S\subseteq U$.

\paragraph{Credit matches degree.} For a matrix-valued $f$ and a vector-valued $g$ of the moves, by
orthogonality and Cauchy--Schwarz,
\begin{gather}
 \mathbb E\bigl[f^\top g\bigr]=\sum_S\mathbb E\bigl[f_S^\top g_S\bigr]=\sum_{r=0}^n\mathbb E\bigl[f_r^\top g_r\bigr],
 \nonumber\\
 \bigl\|\mathbb E[f_r^\top g_r]\bigr\|\le\Bigl(\sum_{|S|=r}\mathbb E\|f_S\|_F^2\Bigr)^{1/2}\Bigl(\sum_{|S|=r}\mathbb E\|g_S\|^2\Bigr)^{1/2}.
 \label{eq:matching}
\end{gather}
At $\hat y=\bm u$ the outcome gradient is $-\pd{\Lout}\theta=\mathbb E[(\Jz)^\top(y-\bm u)]$, so \cref{eq:matching} with $f=\Jz$
and $g=y-\bm u$ is \cref{eq:grad-split}. No permutation property is used.

\begin{lemma}[Gradient split at any parameter value]
\label{lem:gradient-split}
For every differentiable model under the independent-move law, with
$\outcredit(\theta)=\mathbb E\langle\bm z_\theta,y-\bm u\rangle$,
\[
 \dpd{\Lout}{\theta}
 =\underbrace{\mathbb E\bigl[(\Jz)^\top(\hat y-\bm u)\bigr]}_{\text{answer-independent}}
 -\underbrace{\dpd\outcredit\theta}_{\text{answer-dependent}},
\]
\[
 \Bigl\|\dpd\outcredit\theta\Bigr\|\le\sqrt{(A_0+A_1)(J_0+J_1)}+\sqrt{(1-\tfrac1K)\,\Jhigh}.
\]
\end{lemma}

\begin{proof}
Substitute $\hat y-y=(\hat y-\bm u)-(y-\bm u)$ in the per-prompt gradient.
Apply \cref{eq:matching} to the degree groups $\{0,1\}$ and $\{\ge2\}$, using
$\sum_r A_r=\mathbb E\|y-\bm u\|^2=1-1/K$ for the second group.
\end{proof}

\paragraph{Carbery--Wright inequality \citep{carbery2001distributional}.} A polynomial $f$ of degree $d$
under a log-concave law satisfies $\Pr[|f|\le t_{\mathrm{CW}}(\mathbb Ef^2)^{1/2}]\le C_{\mathrm{CW}}\,d\,t_{\mathrm{CW}}^{1/d}$, with $C_{\mathrm{CW}}$ absolute.

\subsection{Target spectrum: unseen-move contraction}
\label{app:target-spectrum}

Every move outside a component's support is averaged out.
The next lemma charges one factor $\mixrate$ for each such move.

\begin{lemma}[Contraction of the target components]
\label{lem:target-contraction}
Assume (A3), and fix the start state and the query. For every $S\subseteq[n]$,
the target has the product representation \cref{eq:product-rep} and satisfies
\[
 \mathbb E\|(y-\bm u)_S\|^2\le(1-\tfrac1K)\mixrate^{2(n-|S|)},
 \qquad A_r\le\tbinom nr(1-\tfrac1K)\mixrate^{2(n-r)}.
\]
\end{lemma}


\begin{proof}
Expand each move into its mean and fluctuation:
\[
 y^\top=\bm e_{s_0}^\top\prod_{t=1}^n
 [\overline{\bm T}+(\bm T_{x_t}-\overline{\bm T})],\qquad
 \bm P_t^{(S)}=
 \begin{cases}\bm T_{x_t}-\overline{\bm T},&t\in S,\\
 \overline{\bm T},&t\notin S.\end{cases}
\]
Each $S$-term is centred in every $x_t$, $t\in S$, and depends only on $x_S$. Hence
\begin{equation}
 (y-\bm u)_S^\top=(\bm e_{s_0}-\bm u)^\top\prod_{t=1}^n\bm P_t^{(S)}.
 \label{eq:product-rep}
\end{equation}
Set $\bm v_0=\bm e_{s_0}-\bm u$ and $\bm v_t^\top=\bm v_{t-1}^\top\bm P_t^{(S)}$.
Every factor preserves $\one^\perp$, so
\begin{align*}
 \|\bm v_t\|^2&\le\mixrate^2\|\bm v_{t-1}\|^2&& (t\notin S),\\
 \E_{x_t}[\|\bm v_t\|^2\mid\bm v_{t-1}]
 &=\|\bm v_{t-1}\|^2-\|\bm v_{t-1}^\top\overline{\bm T}\|^2
 \le\|\bm v_{t-1}\|^2&& (t\in S),\\
 \E\|(y-\bm u)_S\|^2&\le(1-1/K)\mixrate^{2(n-|S|)},\\
 A_r&=\sum_{|S|=r}\E\|(y-\bm u)_S\|^2
 \le\binom nr(1-1/K)\mixrate^{2(n-r)}.
\end{align*}
\end{proof}


\begin{remark}[Counting, and tightness of the exponent]
\label{rem:tight}
\label{cor:counting}
For counting with a query letter of frequency $p$, $\overline{\bm T}=(1-p)\bm I+p\bm P_{\mathrm{cyc}}$ with $\bm P_{\mathrm{cyc}}$ the cyclic shift.
In the Fourier basis $(\bm f_j)_a=\zeta^{jk}/\sqrt K$, $\zeta=e^{2\pi\mathrm i/K}$, every move is diagonal, and with independent
factors in \cref{eq:product-rep},
\begin{gather*}
 \mu_j=1-p+p\zeta^j,
 \qquad
 \mixrate=\max_{j\ne0}|\mu_j|=\sqrt{1-2p(1-p)\bigl(1-\cos\tfrac{2\pi}K\bigr)},
 \\
 A_r=\tbinom nr\tfrac1K\textstyle\sum_{j\ne0}|\mu_j|^{2(n-r)}(1-|\mu_j|^2)^r .
\end{gather*}
The exponent $n-r$ of \cref{lem:target-contraction} cannot be improved. For $X_t=\zeta^{b_t}$, with mean $\nu$ of modulus $\mixrate$, take the
pure degree-$r$ feature and the model
\[
 \psi_S=e^{\mathrm i\alpha}\zeta^{s_0}\frac{\prod_{t\in S}(X_t-\nu)}{(1-\mixrate^2)^{r/2}},
 \quad e^{\mathrm i\alpha}=(\nu/|\nu|)^{n-r},
 \qquad
 z_a=\theta\operatorname{Re}(\zeta^{-a}\psi_S).
\]
Since $\sum_a\zeta^{-a}=0$, $\mathbb E\overline{X_t}=\overline\nu$ for $t\notin S$ and $\mathbb E[\overline{X_t}(X_t-\nu)]=1-\mixrate^2$ for $t\in S$,
\[
 \Bigl|\dpd{\Lout}\theta\Big|_{\theta=0}\Bigr|=\mixrate^{n-r}(1-\mixrate^2)^{r/2},
 \qquad \mathbb E|\psi_S|^2=1 .
\]
The Cauchy--Schwarz bound $\sqrt{A_rJ_r}$ is attained by $\bm z_\theta=\theta(y-\bm u)_r$ at $\theta=0$.
\end{remark}


\begin{remark}[Control: first occurrence]
\label{rem:control}
\label{cor:controls}
First occurrence is not a permutation task, and its answer keeps a degree-$1$ piece that does not
decay. With query frequency $p$, $\mathbf 1[y=1]=b_1$, so for every $n$
\[
 \mathbb E\|(y-\bm u)_{\{1\}}\|^2\ge\mathbb E(b_1-p)^2=p(1-p).
\]
The theory therefore predicts that outcome training can learn it, and it does (\cref{tab:main}).
\end{remark}

%-----------------------------------------------------------------------------

\subsection{Model spectrum: score count, size and support}
\label{app:model-spectrum}

The statement below is the model ingredient used by both credit theorems.
A score adds at most two sites and costs four powers of $\initscale$.
Its technical proof is in \cref{app:network}.

\paragraph{The score expansion.} For $\xi\in[0,1]$ let $\bm z(\xi)$ be the network with every score $\psi_{il}$, in both
blocks, replaced by $\xi\psi_{il}$. Then $\bm z(1)=\bm z_\theta$, and by Taylor's theorem in $\xi$, for every $\bar k\ge0$,
\begin{equation}
 \bm z_\theta=\sum_{k=0}^{\bar k}\bm z^{[k]}+\bm z^{{\mathrm{rem}},\bar k},
 \qquad
 \bm z^{[k]}:=\frac{\bm z^{(k)}(0)}{k!},
 \qquad
 \bm z^{{\mathrm{rem}},\bar k}:=\int_0^1\frac{(1-\xi)^{\bar k}}{\bar k!}\,\bm z^{(\bar k+1)}(\xi)\,d\xi .
 \label{eq:score-split}
\end{equation}

\begin{proposition}[Score count and degree]
\label{prop:score-degree}
Fix the start state and the query, so the sites of the answer position are the letters. For $\theta\in \initball$,
and every $\bar k\ge0$, there is a constant $\attconst$ such that the Jacobian splits as $\Jz=\sum_{k=0}^{\bar k}J^{[k]}+J^{{\mathrm{rem}},\bar k}$ along \cref{eq:score-split} and satisfies
\begin{equation*}
 \begin{gathered}
 \sup\|J^{[0]}\|_F\le C\initscale,
 \qquad
 \sum_{|S|=1}\sup\|J^{[0]}_S\|_F\le C\initscale^2,
 \qquad
 J^{[0]}_S=0\ \ (|S|\ge2),
 \\
 \sum_S\sup\|J^{[k]}_S\|_F\le C\initscale^2(\attconst\initscale^4)^k\ \ (k\ge1,\text{ all sets }S),
 \qquad
 J^{[k]}_S=0\ \ (|S|>2k+1),
 \\
 \sup\|J^{{\mathrm{rem}},\bar k}\|_F\le C\initscale^2(\attconst\initscale^4)^{\bar k+1},
 \end{gathered}
 \tag{P1}\label{eq:P1}
\end{equation*}
and the same holds for $\bm z^{[k]},\bm z^{{\mathrm{rem}},\bar k}$ with one more factor $\initscale$. The suprema are over inputs. The constants do not depend on $n$, $k$ or $\bar k$.
\end{proposition}

Fix $\attconst$ large enough for this proposition and the position-selection bound
\cref{lem:dilution}. It depends on widths and the scaled initialization, not on $n$.
All theorem-level bounds below use this one constant.

\subsection{Spectral overlap: outcome starvation and plateau}
\label{app:thm1}

We prove the parts of \cref{thm:main} in order.

\paragraph{Spectral overlap.} Set $\rho_\initscale=\attconst\initscale^4$ and $q_{\mathrm{att}}=\rho_\initscale/\mixrate^2$, and let
$\operatorname{credit}_k=\|\E[(J^{[k]})^\top(y-\bm u)]\|$ be the credit of the $k$-score part of the Jacobian. For $0<\mixrate<1$,
\cref{lem:target-contraction,prop:score-degree} give
\begin{equation}
 \operatorname{credit}_k
 \le C\initscale^2\rho_\initscale^k\,\mixrate^{\max(n-2k-1,0)}
 \le C\attconst\initscale^6\mixrate^{n-3}q_{\mathrm{att}}^{k-1},\qquad k\ge1:
 \label{eq:score-credit}
\end{equation}
the first factor is the model size and the second the target contraction. When $q_{\mathrm{att}}\le1/2$ these
form a summable geometric bound; \cref{lem:score-sum} also controls the remainder.

\begin{lemma}[Credit summation]
\label{lem:score-sum}
Assume (A1)--(A3), $\theta\in \initball$, $0<\initscale\le\initscale_0$, and $0<\mixrate<1$.
Fix the start state and the query, and set $\bar k=\lceil(n-1)/2\rceil$ and
$q_{\mathrm{att}}=\attconst\initscale^4/\mixrate^2$, where $\attconst$ is as in \cref{prop:score-degree}.
If $q_{\mathrm{att}}\le\frac12$, then
\[
 \biggl\|\sum_{r\ge2}\mathbb E[(\Jz)_r^\top(y-\bm u)_r]\biggr\|
 \le C\initscale^6\mixrate^{n-3}.
\]
Moreover, if $\initscale\le c\mixrate$ with $c\le\min(1,(2\attconst)^{-1/4})$, then
$q_{\mathrm{att}}\le\frac12$ automatically, and $\outcredit(\theta)=\mathbb E\langle\bm z_\theta,y-\bm u\rangle$
satisfies
\begin{equation}
 \Bigl\|\dpd\outcredit\theta\Bigr\|\le C\initscale\mixrate^n,
 \qquad
 |\outcredit|\le C\initscale^2\mixrate^n.
 \label{eq:thm1-full}
\end{equation}
These bounds also hold after averaging over the start state and query.
\end{lemma}


\begin{proof}
Use \cref{lem:target-contraction,prop:score-degree} with $\bar k=\lceil(n-1)/2\rceil$.
For any selected degrees, let $\mathrm{cr}_k$ be the $k$-score credit and $\mathrm{cr}_{\mathrm{rem}}$ the remainder credit.
\begin{align*}
 \|\E[f_S^\top(y-\bm u)_S]\|
 &\le\sup\|f_S\|_F\mixrate^{n-|S|},\\
 \mathrm{cr}_0&\le C\initscale\mixrate^n+C\initscale^2\mixrate^{n-1},\\
 \mathrm{cr}_k&\le C\initscale^2(\attconst\initscale^4)^k\mixrate^{\max(n-2k-1,0)}
 \le C\attconst\initscale^6\mixrate^{n-3}q_{\mathrm{att}}^{k-1},\\
 \mathrm{cr}_{\mathrm{rem}}&\le C\initscale^2(\attconst\initscale^4)^{\bar k+1}
 =C\attconst\initscale^6\mixrate^{2\bar k}q_{\mathrm{att}}^{\bar k}
 \le C\attconst\initscale^6\mixrate^{n-3}q_{\mathrm{att}}^{\bar k}.
\end{align*}
The $k$-score bound uses $\max(n-2k-1,0)\ge n-2k-1$ and $\mixrate<1$.
At degrees $r\ge2$, $\mathrm{cr}_0=0$, so
\[
 \left\|\sum_{r\ge2}\E[(\Jz)_r^\top(y-\bm u)_r]\right\|
 \le C\attconst\initscale^6\mixrate^{n-3}
       \left(\sum_{k=1}^{\bar k}q_{\mathrm{att}}^{k-1}+q_{\mathrm{att}}^{\bar k}\right)
 \le C\initscale^6\mixrate^{n-3}.
\]
At all degrees, including the extra factor $\initscale$ for logits,
\begin{align*}
 \|\nabla\outcredit\|&\le C(\initscale\mixrate^n+\initscale^2\mixrate^{n-1}+\initscale^6\mixrate^{n-3}),\\
 |\outcredit|&\le C(\initscale^2\mixrate^n+\initscale^3\mixrate^{n-1}+\initscale^7\mixrate^{n-3}).
\end{align*}
If $\initscale\le c\mixrate$ with $c\le\min\{1,(2\attconst)^{-1/4}\}$, then
\[
 q_{\mathrm{att}}\le\tfrac12,\qquad
 \frac{\initscale^2\mixrate^{n-1}}{\initscale\mixrate^n}=\frac\initscale\mixrate\le1,\qquad
 \frac{\initscale^6\mixrate^{n-3}}{\initscale\mixrate^n}
 =\left(\frac\initscale\mixrate\right)^5\mixrate^2\le1.
\]
This proves \cref{eq:thm1-full}. Averaging the uniform conditional bounds preserves them.
\end{proof}


\begin{proof}[Proof of \cref{thm:main}]
The hypothesis $\initscale^2\le c\mixrate$ with $c\le(2\attconst)^{-1/2}$ implies
$q_{\mathrm{att}}\le\frac12$. Apply \cref{lem:score-sum}.
\end{proof}


The main-text plateau with a nonzero readout is the following statement.

\begin{proposition}[Plateau with a nonzero readout]
\label{prop:plateau}
Suppose (A1)--(A4) hold with $0<\mixrate<1$ and $\initscale\le c\mixrate$. Gradient descent on $\Lout$ with step size $\stepsize$
satisfies $\theta_j\in\initball$ and $\log K-\Lout(\theta_j)\le C\initscale^2\mixrate^n$ for every $j$ with
$j\stepsize\le c/(\initscale^2+\mixrate^n)$, that is, for $\Omega(\min(\initscale^{-2},\mixrate^{-n})/\stepsize)$ updates.
\end{proposition}

\begin{proof}[Proof of \cref{prop:plateau}]
On $\initball$, \cref{lem:gradient-split,lem:size,lem:score-sum} give
\begin{align*}
 \nabla \Lout&=\E[(\Jz)^\top(\hat y-\bm u)]-\nabla\outcredit,\\
 \sup\|\Jz\|_F&\le C\initscale,\qquad
 \sup\|\hat y-\bm u\|\le C\sup\|\bm z\|\le C\initscale^2.
\end{align*}
Thus
\begin{equation}
 \|\nabla \Lout\|\le C\initscale(\initscale^2+\mixrate^n).
 \label{eq:outcome-grad}
\end{equation}
If the preceding iterates lie in $\initball$, then
\[
 \|\theta_j-\theta_0\|
 \le\stepsize\sum_{i<j}\|\nabla \Lout(\theta_i)\|
 \le Cj\stepsize\initscale(\initscale^2+\mixrate^n)\le\ballradius\initscale
 \quad\left(j\stepsize\le\frac{\ballradius/C}{\initscale^2+\mixrate^n}\right).
\]
Induction keeps every such iterate in $\initball$. Finally,
\begin{align*}
 \log\sum_ae^{z_a}&\ge\log K+\langle\bm u,\bm z\rangle,\\
 \log K-\Lout(\theta_j)&\le\E\langle\bm z_{\theta_j},y-\bm u\rangle
 =\outcredit(\theta_j)\le C\initscale^2\mixrate^n,\\
 \frac1{\initscale^2+\mixrate^n}&\ge\tfrac12\min\{\initscale^{-2},\mixrate^{-n}\}.
\end{align*}
\end{proof}


\begin{remark}[Credit along a direction]
\label{rem:direction}
For any direction $\bm d$, \cref{eq:matching} with $f=(\Jz)\bm d$ gives
$\langle\pd\outcredit\theta,\bm d\rangle=\sum_r\mathbb E\langle((\Jz)\bm d)_r,(y-\bm u)_r\rangle$. Each credit norm bound applies to every unit direction.
For general $\bm d$, multiply it by $\|\bm d\|$.
This includes directions toward an exact outcome solution.
\end{remark}

\begin{remark}[Where the cap $\initscale^{-2}$ comes from]
\label{rem:cap}
Initial logits put predictions $O(\initscale^2)$ away from uniform, which gives the cap. With $\tilde V=0$
the cap disappears (\cref{thm:outcome-end}). The runs of \cref{fig:regime} use a zero readout; the other experiments
use a nonzero one.
\end{remark}

\begin{remark}[The case $\mixrate=0$]
\label{rem:lambda-zero}
If $\mixrate=0$, the hypotheses of \cref{thm:main} force $\initscale=0$, so this case needs its own statement.
\Cref{lem:target-contraction} gives $(y-\bm u)_S=0$ unless $S=[n]$. Take $\bar k=\lceil(n-1)/2\rceil$ in \cref{prop:score-degree}. The parts $J^{[k]}$ with
$2k+1<n$, that is $k<\bar k$, have no component on $[n]$, so only $J^{[\bar k]}$ and $J^{{\mathrm{rem}},\bar k}$ pair with the answer.
For $\initscale\le\initscale_0$ with $\attconst\initscale_0^4\le1$, the remainder is at most the $\bar k$-score term, and
\[
 \Bigl\|\dpd\outcredit\theta\Bigr\|\le C\initscale^2(\attconst\initscale^4)^{\bar k},
 \qquad
 |\outcredit|\le C\initscale^3(\attconst\initscale^4)^{\bar k},
 \qquad
 4\bar k+2=\begin{cases}2n,&n\text{ odd},\\2n+2,&n\text{ even}.\end{cases}
\]
The proof of \cref{prop:plateau} holds with $\mixrate^n$ replaced by $0$: the loss stays within $C\initscale^3(\attconst\initscale^4)^{\bar k}$ of chance for a training
time of order $\initscale^{-2}$. So the answer-dependent gradient is of order $\initscale^{4\bar k+2}$, an exponent of at least $2n$, where the
second-order expansion gave only $\initscale^{10}$.
Counting modulo $2$ at match probability $\frac12$ is this case, and so is BIG-Bench Hard tracking
with three objects. With a zero readout, the proof of \cref{thm:outcome-end} holds with
$C_2\initscale^2\mixrate^n$ replaced by $C\initscale^3(\attconst\initscale^4)^{\bar k}$, and the induction in its proof allows
$j\stepsize\le c\,\initscale^{-1}(\attconst\initscale^4)^{-\bar k}$, a time of order at least $\initscale^{-(2n-1)}$.
\label{rem:end-parity}
\end{remark}

\subsection{Local targets: step rule credit}
\label{app:thm2}
\label{app:separation}
\label{app:rule}

We first identify contributing terms, then bound a nonzero readout derivative.

\paragraph{Notation.} Let $\bm g_t=\bm e_{s_t}-\overline{\bm T}^\top\bm e_{s_{t-1}}$ and $\rulecredit(\theta)=\mathbb E_t\mathbb E\langle\bm z_t,\bm g_t\rangle$, so $\rulecred:=\|\partial\rulecredit/\partial\theta\|$. Average over $t=2,\ldots,n$.
Letter $w_t$ sits at $l_t=1+t$.
Position $i_t=n+1+t$ predicts $s_t$ and holds $s_{t-1}$.
Thus $m_{l_t}=t+2$ and $m_{i_t}=n+2+t$. The tokens visible from $i_t$ are $s_0$, the query, the letters and the displayed
states, each a function of $(s_0,q,w_{<t})$. \Cref{lem:size,lem:support,lem:dilution} hold at every position.
\begin{lemma}[Average positional dilution]
\label{lem:dilution-average}
For $n\ge2$ and the step positions above, with $t$ uniform on $\{2,\ldots,n\}$,
\begin{equation}
 \mathbb E_t\,\kappa_{l_t,i_t}\le\frac1{n+4}\Bigl(1+\log2+\frac{n+2}{n-1}\Bigr)\le\frac6n .
 \label{eq:dilution-avg}
\end{equation}
\end{lemma}


\begin{proof}
Using $m_j=j+1$, $m_{i_t}\ge n+4$ and an integral bound,
\begin{align*}
 \sum_{l_t<j\le i_t}\frac1{m_j}
 &\le\log2+\log\frac{n+2}{t+2},\\
 \sum_{t=2}^n\log\frac{n+2}{t+2}
 &\le\sum_{k=1}^{n+2}\log\frac{n+2}{k}\le n+2,\\
 \E_t\kappa_{l_t,i_t}
 &\le\frac1{n+4}\left(1+\log2+\frac{n+2}{n-1}\right)\le\frac6n.
\end{align*}
\end{proof}


\begin{lemma}[Only the letter and a state earn rule credit]
\label{lem:rule-sites}
Under (A3$'$), a term $\term$ of $\bm z_t$ or $h^2_{i_t}$ has $\mathbb E\langle \term,\bm g_t\rangle=0$ unless it contains the letter input $x_{l_t}$ and
the input of a state ($s_0$ or a displayed state). In particular no term of $\bm z^{[0]}_t$ contributes.
\end{lemma}


\begin{proof}
Let $A$ be the term's input sites. If $l_t\notin A$, condition on $(s_0,q,w_{\ne t})$:
\[
 \E[\bm g_t\mid s_0,q,w_{\ne t}]=0,
 \qquad\E\langle\term,\bm g_t\rangle=0.
\]
If $A$ contains no state input, the term is a function of $(q,w)$, and (A3$'$) gives
\[
 \E[\bm g_t\mid q,w]=\bm u-\overline{\bm T}^\top\bm u=0,
 \qquad\E\langle\term,\bm g_t\rangle=0.
\]
Every contributing term has at least two sites; \cref{lem:support}(a) excludes $k=0$.
\end{proof}


For readout indices $a,b$, let
\[
 Q_t:=\sum_{\ell=1,2}\ \sum_{\mathrm{heads}}\ \frac1{m_{i_t}}\sum_{j\le i_t}W_Vx_j\bigl(\psi^{(0)}_{i_tj}-\bar\psi^{(0)}_{i_t}\bigr),
\]
\[
 \psi^{(0)}_{il}=\frac{(W_Qx_i)^\top W_Kx_l}{\sqrt{d_{\mathrm{h}}}},
 \qquad
 \bar\psi^{(0)}_i=\frac1{m_i}\sum_{l\le i}\psi^{(0)}_{il},
\]
the one-score term of \cref{lem:support}(c) with every residual replaced by its input, and
$P_n(\tilde\theta):=\mathbb E_t\,\mathbb E\bigl[(\bm g_t)_a(Q_t)_b\bigr]\big|_{\theta=\tilde\theta}$.

\begin{lemma}[Leading rule coefficient]
\label{lem:leading}
Assume (A3$'$). Fix a readout entry $V_{ab}$ and define the directional derivative
\[
 F_n(\initscale):=\left.\frac{\partial\rulecredit(\theta)}{\partial V_{ab}}\right|_{\theta=\initscale\tilde\theta}
 =\left.\mathbb E_t\mathbb E[(\bm g_t)_a h^2_{i_t,b}]\right|_{\theta=\initscale\tilde\theta}.
\]
The coordinate direction has unit norm, hence $\rulecred(\theta_0)\ge|F_n(\initscale)|$.
Let $P_n$ be defined from $Q_t$ above.
\begin{enumerate}[label=(\roman*),leftmargin=*,itemsep=0pt,topsep=2pt]
 \item $P_n$ is a homogeneous polynomial of degree $6$ in the attention and embedding entries of $\tilde\theta$, the
 same for every activation pattern, and for $\initscale\le\initscale_0$
 \[
 F_n(\initscale)=\initscale^6P_n(\tilde\theta)+\varrho_n(\initscale),
 \qquad
 |\varrho_n(\initscale)|\le\frac{C'}{\initscale_0}\cdot\frac{\initscale^7}n+\frac{C\initscale^{10}}n .
 \]
 \item For counting with a fixed query of match probability $p<\frac12$, $P_n$ contains a monomial $\prod_{y\in \mathcal E}y$ in six
 entries of $\tilde\theta$ with coefficient $a_{\mathcal E}\ge a_{\mathrm{match}}/n$ for every $n\ge2$, where $a_{\mathrm{match}}:=\frac{p(1-p)\min(1-2p,\frac14)}{3K\sqrt{d_{\mathrm{h}}}}>0$.
\end{enumerate}
\end{lemma}


\begin{proof}
\proofstep{(i) Leading term.}
Take $\bar k=2+\lceil\log_2 n\rceil$ and $\rho_\initscale=\attconst\initscale^4\le1/2$.
By \cref{lem:rule-sites,lem:dilution,eq:dilution-avg},
\begin{align*}
 \E[(\bm g_t)_a h_{i_t,b}^{2,[0]}]&=0,\\
 2C\initscale^2\E_t\kappa_{l_t,i_t}\sum_{k=2}^{\bar k}\rho_\initscale^k
 &\le C\initscale^{10}/n,\\
 2C\initscale^2\rho_\initscale^{\bar k+1}
 &\le C\attconst^2\initscale^{10}2^{-(\bar k-1)}\le C\initscale^{10}/n.
\end{align*}
With fixed activation matrices, one-score monomials satisfy
\[
 \term(\initscale\tilde\theta)=\initscale^{\deg\term}\term(\tilde\theta),\qquad
 \deg\term\ge6.
\]
Degree $6$ consists exactly of $Q_t$. An MLP matrix correction adds two degrees;
a bias output replacing an input adds one. The higher-degree terms obey
\[
 \initscale^{\deg\term}\le(\initscale/\initscale_0)^7\initscale_0^{\deg\term},\qquad
 2(\initscale/\initscale_0)^7C\initscale_0^6\E_t\kappa_{l_t,i_t}
 \le\frac{C'}{\initscale_0}\frac{\initscale^7}{n}.
\]
Consequently,
\[
 F_n(\initscale)=\initscale^6P_n(\tilde\theta)+\varrho_n(\initscale),\qquad
 |\varrho_n(\initscale)|\le\frac{C'}{\initscale_0}\frac{\initscale^7}{n}+\frac{C\initscale^{10}}n.
\]

\proofstep{(ii) A nonzero coefficient.}
Choose a block-2 head, a state $s^\ast$, an output coordinate $b$ and $a=s^\ast+1$.
Retain only these six entries:
\[
 \mathcal E=\{\tilde E_{s^\ast,0},\tilde E_{q,1},\tilde E_{q,2},
       (\tilde W_Q)_{0,0},(\tilde W_K)_{0,1},(\tilde W_V)_{b,2}\}.
\]
Set them to $1$ and all others to $0$, giving $\tilde\theta^{\mathcal E}$.
The sole surviving monomial is $\prod_{y\in \mathcal E}y$, with coefficient $a_{\mathcal E}=P_n(\tilde\theta^{\mathcal E})$.
Writing $\phi_j=[\text{token at }j=q]$ and $m=n+2+t$,
\begin{align*}
 N&=\sum_{j\le i_t}\phi_j=1+b_t+N',\\
 W_Vx_j&=\phi_j\bm e_b,\qquad
 \psi^{(0)}_{i_tj}=\frac{[s_{t-1}=s^\ast]\phi_j}{\sqrt{d_{\mathrm{h}}}},\\
 (Q_t)_b&=\frac{[s_{t-1}=s^\ast]}{\sqrt{d_{\mathrm{h}}}}
                  \left(\frac Nm-\frac{N^2}{m^2}\right),\\
 \bm g_t&=(b_t-p)(\bm e_{s_{t-1}+1}-\bm e_{s_{t-1}}),\\
 \E[(b_t-p)N]&=p(1-p),\\
 \E[(b_t-p)N^2]&=p(1-p)[3+2(n-1)p].
\end{align*}
The uniform start makes $s_{t-1}$ independent of all letters. Therefore
\begin{align*}
 a_{\mathcal E}&=\frac1{K\sqrt{d_{\mathrm{h}}}(n-1)}
      \sum_{t=2}^n\frac{p(1-p)}{n+2+t}
       \left(1-\frac{3+2(n-1)p}{n+2+t}\right),\\
 3+2(n-1)p&\le\max\{2p,3/4\}(n+4),\\
 a_{\mathcal E}&\ge\frac{p(1-p)\min\{1-2p,1/4\}}{K\sqrt{d_{\mathrm{h}}}(2n+2)}
 \ge\frac{a_{\mathrm{match}}}n.
\end{align*}
\end{proof}

\begin{proof}[Proof of \cref{thm:separation}(i)] At the position $i$ predicting $s_t$,
$h^{(0)}_i=E_{\mathrm{val}}^\top\bm e_{s_{t-1}}+P\bm e_i$, and \cref{lem:size}(a) gives $h_i=h^{(0)}_i+O(\initscale^2)$ and $\hat y_i-\bm u=O(\initscale^2)$.
Using $\mathbb E[\bm e_{s_t}-\bm u]=0$ and $\mathbb E[\bm e_{s_t}\bm e_{s_{t-1}}^\top]=\frac1K\overline{\bm T}^\top$,
\[
 \dpd{\Lstep^{\ge2}}{V}=-\mathbb E\bigl[(\bm e_{s_t}-\bm u)h_i^{(0)\top}\bigr]+O(\initscale^2)=-\bm\Gamma E_{\mathrm{val}}+O(\initscale^2),
 \qquad
 \bm\Gamma=\tfrac1K\bigl(\overline{\bm T}^\top-\tfrac1K\one\one^\top\bigr).
\]
By (A2), $E_{\mathrm{val}}$ has full row rank and
$\varsigma_{\min}(E_{\mathrm{val}})\ge\initscale\varsigma/2$ on $\initball$.
The largest singular value of $\bm\Gamma$ is $\mixrate/K$.
For $\initscale\le\mixrate\varsigma/(4KC)$, this gives
\[
 \Bigl\|\dpd{\Lstep^{\ge2}}{\theta}\Bigr\|\ge\Bigl\|\dpd{\Lstep^{\ge2}}{V}\Bigr\|_F\ge\|\bm\Gamma E_{\mathrm{val}}\|_F-C\initscale^2\ge\frac{\initscale\mixrate\varsigma}{4K}.
\]
The outcome bound is \cref{eq:outcome-grad}.
\end{proof}


\begin{proof}[Proof of \cref{thm:separation}(ii)]
Set $\rho_\initscale=\attconst\initscale^4\le1/2$ and $\bar k=1+\lceil\log_2 n\rceil$.
By \cref{lem:rule-sites,lem:dilution,lem:dilution-average},
\begin{align*}
 \rulecred(\theta)
 &\le2C\initscale^2\E_t\kappa_{l_t,i_t}\sum_{k=1}^{\bar k}\rho_\initscale^k
     +2C\initscale^2\rho_\initscale^{\bar k+1}\\
 &\le C\attconst\initscale^6\left(\E_t\kappa_{l_t,i_t}+2^{-\bar k}\right)
 \le C\attconst\initscale^6\left(\frac6n+\frac1{2n}\right)
 \le\frac{C\initscale^6}{n}.
\end{align*}
The remainder uses \cref{lem:size}(b), with constants independent of $\bar k$.
\end{proof}


\begin{proof}[Proof of \cref{thm:separation}(iii)]
Fix $n$. The non-bias entries are independent uniform variables on $[-1,1]$.
The monomial in \cref{lem:leading}(ii) is the unique degree-$6$ term with support $\mathcal E$, so
\[
 \E P_n(\tilde\theta)^2\ge a_{\mathcal E}^2\prod_{y\in \mathcal E}\E y^2
 =a_{\mathcal E}^2\,3^{-6}\ge\frac{a_{\mathrm{match}}^2}{729n^2}.
\]
For $t_{\mathrm{CW}}=(\failprob/(6C_{\mathrm{CW}}))^6$, Carbery--Wright gives
\[
 \Pr\left(|P_n(\tilde\theta)|\ge\frac{t_{\mathrm{CW}} a_{\mathrm{match}}}{27n}\right)\ge1-\failprob.
\]
The bounded entries and fixed biases give common constants in \cref{lem:leading}(i).
Choose
\[
 \initscale_\failprob=\min\left\{\initscale_0,
    \frac{t_{\mathrm{CW}} a_{\mathrm{match}}\initscale_0}{108C'},
    \left(\frac{t_{\mathrm{CW}} a_{\mathrm{match}}}{108C}\right)^{1/4}\right\}.
\]
On the preceding event, for every $\initscale\le\initscale_\failprob$,
\begin{align*}
 |\varrho_n(\initscale)|
 &\le\frac{C'}{\initscale_0}\frac{\initscale^7}n+\frac{C\initscale^{10}}n
 \le\frac{t_{\mathrm{CW}} a_{\mathrm{match}}}{54}\frac{\initscale^6}n,\\
 \rulecred(\theta_0)&\ge|F_n(\initscale)|
 \ge\initscale^6|P_n|-|\varrho_n(\initscale)|
 \ge\frac{t_{\mathrm{CW}} a_{\mathrm{match}}}{54}\frac{\initscale^6}n
 =a_\failprob\frac{\initscale^6}n.
\end{align*}
\end{proof}


%-----------------------------------------------------------------------------
\subsection{End of training: basins and the outcome plateau}
\label{app:end}

This subsection proves \cref{thm:outcome-end,thm:basin,cor:exact}. The table lists what each result uses.

\begin{center}\small
\fitwidth{\begin{tabular}{@{}lll@{}}
\toprule
Result & Uses & Hypotheses\\
\midrule
\Cref{lem:end-barrier}: small loss forces margins & definitions only & $L(\theta)<\pimin/4$\\
\Cref{lem:end-readout}: readout direction & \cref{lem:end-barrier} & $L(\theta)<\pimin/4$, logits $Vh$\\
\Cref{thm:basin}: convergence in the basin & \cref{lem:end-barrier,lem:end-readout} & \cref{eq:descent} for $j\ge \entryidx$, $\theta_{\entryidx}\in\basin(L)$\\
\Cref{lem:end-nonempty}: the basin contains the executor & proof of \cref{prop:represent} & positive margins\\
\Cref{cor:exact}: exact generation & \cref{thm:basin} & as \cref{thm:basin}, for $\Lstep^{\mathrm{full}}$\\
\Cref{lem:end-smooth}: smoothness near $\theta_0$ & \cref{lem:softmax-all}, fixed ReLU pattern & (A1), (A2) on $\initball'$\\
\Cref{thm:outcome-end}: outcome does not converge & \cref{lem:score-sum,lem:end-smooth}, Pinsker & (A1)--(A3), $\initscale\le c\mixrate$, $\tilde V=0$\\
\bottomrule
\end{tabular}}
\end{center}

\paragraph{Setting.} $L=\sum_{i\in\items}\pi_i\ell_i$ is a supervised loss as in \cref{sec:end}.
For item $i$ with logits $\bm z_i=Vh_i$ and gold token $y_i$, write
$\hat y_i=\softmax(\bm z_i)$, $q_i=\hat y_{i,y_i}$ and
$\margin_i=\min_{a\ne y_i}(z_{i,y_i}-z_{i,a})$.

\paragraph{Small loss forces correct predictions.}

\begin{lemma}[Correctness barrier]
\label{lem:end-barrier}
If $L(\theta)<\pimin/4$, then for every item $i$: $\ell_i<1/4$, $q_i>1/2$, the gold token has the
strictly largest logit, and
\[
 \margin_i\ \ge\ \log\frac{\pi_i}{2L(\theta)}\ \ge\ \log\frac{\pimin}{2L(\theta)}\ >\ \log2 .
\]
\end{lemma}

\begin{proof}
Write $L=L(\theta)$. Every term of $L$ is nonnegative, so $\pi_i\ell_i\le L<\pimin/4\le\pi_i/4$; hence $\ell_i<1/4$
and $q_i=e^{-\ell_i}>e^{-1/4}>1/2$. The sum in
$\ell_i=\log\bigl(1+\sum_{a\ne y_i}e^{-(z_{i,y_i}-z_{i,a})}\bigr)$ contains the term of the competitor that attains
$\margin_i$, so $\ell_i\ge\log(1+e^{-\margin_i})$ and $e^{-\margin_i}\le e^{\ell_i}-1\le2\ell_i$, using $e^x-1\le2x$ on
$[0,1]$. Therefore, with $\ell_i\le L/\pi_i$ and $L<\pimin/4$,
\[
 \margin_i\ \ge\ \log\frac1{2\ell_i}\ \ge\ \log\frac{\pi_i}{2L}\ \ge\ \log\frac{\pimin}{2L}\ >\ \log2 .
\]
In particular $\margin_i>0$, so the gold logit is strictly largest.
\end{proof}

\begin{lemma}[The readout direction decreases the loss]
\label{lem:end-readout}
If $L(\theta)<\pimin/4$ and $\margin=\log(\pimin/(2L(\theta)))$, then
\[
 \langle V,\nabla_VL(\theta)\rangle_F\le-\tfrac\margin2\,L(\theta),
 \qquad
 \|\nabla L(\theta)\|\ \ge\ \|\nabla_VL(\theta)\|_F\ \ge\ \frac{\margin\,L(\theta)}{2\|V\|_F}.
\]
\end{lemma}

\begin{proof}
Since $\bm z_i=Vh_i$ and $h_i$ does not depend on $V$, $\nabla_V\ell_i=(\hat y_i-\bm e_{y_i})h_i^\top$ and
$\langle V,\nabla_V\ell_i\rangle_F=-\langle\bm z_i,\bm e_{y_i}-\hat y_i\rangle$. As $\sum_a\hat y_{i,a}=1$,
\[
 \langle\bm z_i,\bm e_{y_i}-\hat y_i\rangle=\sum_{a\ne y_i}\hat y_{i,a}(z_{i,y_i}-z_{i,a})
 \ \ge\ \margin_i(1-q_i)\ \ge\ \tfrac12\margin_i\ell_i ,
\]
where the last step uses $\ell_i=-\log q_i\le(1-q_i)/q_i\le2(1-q_i)$, valid because $q_i>1/2$ by
\cref{lem:end-barrier}. Multiplying by $\pi_i$, summing, and using $\margin_i\ge\margin$ (\cref{lem:end-barrier}) gives
$\langle V,\nabla_VL\rangle_F\le-\frac\margin2L$. The second claim follows by Cauchy--Schwarz, since $\nabla_VL$ is a block
of $\nabla L$; here $V\ne0$, because $V=0$ would make all logits equal.
\end{proof}

\paragraph{Convergence in the basin.}

\begin{proof}[Proof of \cref{thm:basin}]
Write $L_j=L(\theta_j)$ and $V_j$ for the readout of $\theta_j$. By \cref{eq:descent}, which \cref{ass:armijo} guarantees,
the losses are non-increasing for $j\ge\entryidx$, so $L_j\le L_{\entryidx}<\pimin/4$ and $\theta_j\in\basin(L)$.
\Cref{lem:end-barrier} then gives part~(i), every margin is at least $\log(\pimin/(2L_j))\ge\margin_{\entryidx}>\log2$,
and $V_j\ne0$.

Summing \cref{eq:descent} over $i=\entryidx,\ldots,j-1$ gives $\sum_i\stepsize_i\|\nabla L(\theta_i)\|^2\le2L_{\entryidx}$.
Since $V$ is a block of $\theta$, the triangle and Cauchy--Schwarz inequalities give
$\|V_j-V_{\entryidx}\|_F\le\sum_i\stepsize_i\|\nabla L(\theta_i)\|\le(\traintime_j\sum_i\stepsize_i\|\nabla L(\theta_i)\|^2)^{1/2}
\le(2L_{\entryidx}\traintime_j)^{1/2}$, hence
\begin{equation}
 \|V_j\|_F^2\le2\|V_{\entryidx}\|_F^2+4L_{\entryidx}\traintime_j .
 \label{eq:end-vgrowth}
\end{equation}
At each $i\ge\entryidx$, \cref{eq:descent} and \cref{lem:end-readout}, which applies because $\theta_i\in\basin(L)$, give
\[
 L_{i+1}\le L_i-\frac{\stepsize_i}2\|\nabla L(\theta_i)\|^2\le L_i(1-a_iL_i),
 \qquad a_i:=\frac{\stepsize_i\margin_{\entryidx}^2}{8\|V_i\|_F^2}.
\]
Since $L_{i+1}>0$, $0\le a_iL_i<1$, and $1/(1-x)\ge1+x$ gives $1/L_{i+1}\ge1/L_i+a_i$. Sum over $i$, bound $\|V_i\|_F^2$ by
\cref{eq:end-vgrowth}, and compare the sum with the integral of the decreasing function $\traintime\mapsto1/(a+b\traintime)$,
where $a=2\|V_{\entryidx}\|_F^2$ and $b=4L_{\entryidx}$, using $\traintime_{i+1}=\traintime_i+\stepsize_i$:
\[
 \begin{aligned}
 \frac1{L_j}&\ge\frac1{L_{\entryidx}}+\frac{\margin_{\entryidx}^2}8\sum_{i=\entryidx}^{j-1}\frac{\stepsize_i}{a+b\traintime_i}
 \ge\frac1{L_{\entryidx}}+\frac{\margin_{\entryidx}^2}8\int_0^{\traintime_j}\frac{d\traintime}{a+b\traintime}\\
 &=\frac1{L_{\entryidx}}\Bigl(1+\frac{\margin_{\entryidx}^2}{32}\log\bigl(1+2L_{\entryidx}\traintime_j/\|V_{\entryidx}\|_F^2\bigr)\Bigr).
 \end{aligned}
\]
Inverting gives part~(ii).
\end{proof}

\begin{lemma}[The basin contains the executor]
\label{lem:end-nonempty}
Let $\theta^\star$ have logit margin at least $\execmargin>0$ on every supervised item, and let $\Ktokens$ be
the number of logits in the softmax. Replacing the readout $V^\star$ by $aV^\star$ with
$a>(\execmargin)^{-1}\log(4(\Ktokens-1)/\pimin)$ gives a point of $\basin(L)$, and $\basin(L)$ is open.
For the executors $\thstep$ and $\thout$ of \cref{prop:represent}, $\execmargin>19$ on every
supervised position, so this holds for $\Lstep$, $\Lstep^{\mathrm{full}}$ and $\Lout$.
\end{lemma}

\begin{proof}
Since $h_i$ does not depend on $V$, replacing $V^\star$ by $aV^\star$ multiplies every logit, hence every margin, by $a$.
Each of the $\Ktokens-1$ competitors then trails the gold logit by at least $a\execmargin$, so
$\ell_i\le\log(1+(\Ktokens-1)e^{-a\execmargin})\le(\Ktokens-1)e^{-a\execmargin}<\pimin/4$ by the choice of $a$, and summing
with weights $\pi_i$ gives $L<\pimin/4$. The set $\basin(L)$ is open because $L$ is continuous. For the executors, the
finite-score step of the proof of \cref{prop:represent} gives margin greater than $19$ at every supervised position,
including the formatting and answer positions.
\end{proof}

\begin{proof}[Proof of \cref{cor:exact}]
By \cref{thm:basin}(i) for $\Lstep^{\mathrm{full}}$, at every $j\ge\entryidx$ the gold token has the strictly largest
teacher-forced logit at every generated position of every supported prompt. Induct on the generated positions: if
the first $t$ greedy tokens agree with the gold continuation, the prefix is the gold prefix, so the next greedy token
is the teacher-forced prediction and hence gold. The gold continuation ends in EOS, where greedy decoding stops.
\end{proof}

\paragraph{Enlarged ball.} For \cref{thm:outcome-end}, take $\ballradius\le\min(1/8,\varsigma/4)$ in (A2) and let
$\initball'=\{\theta:\|\theta-\theta_0\|\le2\ballradius\initscale\}$. Then $2\ballradius$ satisfies the constraint of (A2),
so every bound of \cref{app:mainproofs} that holds on $\initball$ also holds on $\initball'$, with constants
independent of $n$, $\initscale$ and $\stepsize$. A zero readout in $\tilde\theta$ is admissible: (A2) constrains the MLP biases
and the value-token embeddings, not the readout, and a zero block only lowers the norms in those bounds.

\begin{lemma}[Smoothness near initialization]
\label{lem:end-smooth}
Assume (A1) and (A2), with $\initscale\le\initscale_0$. There is $C_1$, independent of $n$ and $\initscale$, such that for every
$\theta\in\initball'$ and every direction $\bm v$,
\[
 \|\nabla \Lout(\theta)\|\le C_1\initscale,
 \qquad
 |\partial^2\Lout(\theta)[\bm v,\bm v]|\le C_1\|\bm v\|^2 .
\]
\end{lemma}

\begin{proof}
Fix a unit direction $\bm v$ and write $F'$, $F''$ for derivatives along $\xi\mapsto\theta+\xi\bm v$. Say $F$ has
\emph{profile} $(a,b,c)$ if $\|F\|\le a$, $\|F'\|\le b$ and $\|F''\|\le c$ on $\initball'$ for every input; each $C$ below
is independent of $n$, $\initscale$, the input and the position. Every parameter block has profile $(a_0\initscale,1,0)$,
with $a_0$ as in \cref{app:network}, and every input $x_j=E\bm e_{\mathrm{tok}_j}+P_{\mathrm{pos}}\bm e_j$ has profile
$(a_0\initscale,2,0)$. By the product rule $(FG)''=F''G+2F'G'+FG''$, a bilinear product of norm at most $1$ of factors
with profiles $(a_F,b_F,c_F)$ and $(a_G,b_G,c_G)$ has profile $(a_Fa_G,\,b_Fa_G+a_Fb_G,\,c_Fa_G+2b_Fb_G+a_Fc_G)$.

We show by induction over the blocks that every residual $h^\ell_i$ has profile $(C\initscale,C,C)$; for $\ell=0$ this is
the input bound. The score $\psi_{ij}=(W_Qh_i)^\top W_Kh_j/\sqrt{d_{\mathrm h}}$ is a product of four factors of profile
$(C\initscale,C,C)$, so $|\psi_{ij}|\le C\initscale^4$, $|\psi_{ij}'|\le C\initscale^3$ and $|\psi_{ij}''|\le C\initscale^2$,
since a second derivative leaves at least two factors undifferentiated. By \cref{lem:softmax-all} with $k=1,2$, the
attention row satisfies $\|\alpha_{i\cdot}'\|_1\le2\|\psi_{i\cdot}'\|_\infty\le C\initscale^3$ and
$\|\alpha_{i\cdot}''\|_1\le8\|\psi_{i\cdot}'\|_\infty^2+2\|\psi_{i\cdot}''\|_\infty\le C\initscale^2$ on any number of positions,
and $\|\alpha_{i\cdot}\|_1=1$. A head output $\bm a_i=\sum_{j\le i}\alpha_{ij}\bm u_j$ with $\bm u_j=W_Vh_j$ of profile
$(C\initscale^2,C\initscale,C)$ therefore satisfies
\[
 \|\bm a_i''\|\le\|\alpha_{i\cdot}''\|_1\max_j\|\bm u_j\|+2\|\alpha_{i\cdot}'\|_1\max_j\|\bm u_j'\|+\max_j\|\bm u_j''\|\le C,
\]
and similarly $\|\bm a_i\|\le C\initscale^2$ and $\|\bm a_i'\|\le C\initscale$; only $\ell_1$ norms of attention rows enter, so
the bounds do not depend on the number of positions. On $\initball'$ the ReLU pattern is fixed, so
$h^\ell_i=(h^{\ell-1}_i+\bm a^\ell_i)+(M_\ell-I)(h^{\ell-1}_i+\bm a^\ell_i)+\bm o_\ell$, where $M_\ell-I=W_{\mathrm{out}}DW_{\mathrm{in}}$ and
$\bm o_\ell=W_{\mathrm{out}}D\bm b$ are products of two parameter blocks with profile $(C\initscale^2,C\initscale,C)$. The product
rule gives $h^\ell_i$ profile $(C\initscale,C,C)$, which completes the induction.

The logits $\bm z=Vh^2$ then have profile $(C\initscale^2,C\initscale,C)$. For one prompt, $\ell=\log\sum_ae^{z_a}-z_y$, so
\[
 |\ell'|=|\langle\hat y-\bm e_y,\bm z'\rangle|\le\sqrt2\,C\initscale,
 \qquad
 |\ell''|=\bigl|\bm z'^\top(\operatorname{diag}\hat y-\hat y\hat y^\top)\bm z'+\langle\hat y-\bm e_y,\bm z''\rangle\bigr|
 \le\tfrac12C^2\initscale^2+\sqrt2\,C,
\]
using $\|\hat y-\bm e_y\|\le\sqrt2$ and that the softmax covariance has operator norm at most $1/2$: for a unit $v$,
$v^\top(\operatorname{diag}\hat y-\hat y\hat y^\top)v=\operatorname{Var}_{\hat y}(v)\le\frac14(\max_av_a-\min_av_a)^2\le\frac12$.
Averaging over prompts and taking the supremum over unit $\bm v$ gives the claim with $C_1=\frac12C^2+\sqrt2C$.
\end{proof}

\paragraph{The outcome plateau without a cap.}

\begin{proof}[Proof of \cref{thm:outcome-end}]
Let $C_1$ be the constant of \cref{lem:end-smooth} and $C_2$ that of \cref{lem:score-sum} on $\initball'$, valid for
$\initscale\le c\mixrate$, so that $|\outcredit(\theta)|\le C_2\initscale^2\mixrate^n$ on $\initball'$. Take
$\stepsize\le\stepsize_0:=\min(1,\ballradius)/C_1$ and $c_{\mathrm{time}}:=\ballradius^2/(2C_2)$, and decrease $\initscale_0$
so that $C_2\initscale_0^2\le1/4$.

For one prompt, $\KL(\bm u\|\hat y)=\log\sum_ae^{z_a}-\log K-\langle\bm u,\bm z\rangle$, so
$\ell=\log K+\KL(\bm u\|\hat y)-\langle\bm z,\bm e_y-\bm u\rangle$. Taking expectations,
\begin{equation}
 \Lout(\theta)=\log K+\E\,\KL(\bm u\,\|\,\hat y_\theta)-\outcredit(\theta).
 \label{eq:lo-identity}
\end{equation}
At $\theta_0$ the readout is zero, so $\bm z=0$, $\hat y=\bm u$ and $\Lout(\theta_0)=\log K$. Since $\KL\ge0$,
\cref{eq:lo-identity} gives the floor
\begin{equation}
 \Lout(\theta)\ge\log K-C_2\initscale^2\mixrate^n,\qquad\theta\in\initball'.
 \label{eq:lo-floor}
\end{equation}

We show by induction that $\theta_j\in\initball$ whenever $j\stepsize\le c_{\mathrm{time}}\mixrate^{-n}$. Suppose
$\theta_0,\ldots,\theta_j\in\initball$ and $(j+1)\stepsize\le c_{\mathrm{time}}\mixrate^{-n}$. For $i\le j$, \cref{lem:end-smooth}
gives $\|\theta_{i+1}-\theta_i\|=\stepsize\|\nabla\Lout(\theta_i)\|\le\stepsize_0C_1\initscale\le\ballradius\initscale$, so the
segment from $\theta_i$ to $\theta_{i+1}$ lies in $\initball'$, where $\Lout$ is smooth with Hessian bounded by $C_1$.
Taylor's theorem with integral remainder and $C_1\stepsize\le1$ give
\[
 \Lout(\theta_{i+1})\le\Lout(\theta_i)-\stepsize\|\nabla\Lout(\theta_i)\|^2\bigl(1-\tfrac{C_1\stepsize}2\bigr)
 \le\Lout(\theta_i)-\tfrac\stepsize2\|\nabla\Lout(\theta_i)\|^2 .
\]
Summing over $i\le j$ and applying \cref{eq:lo-floor} at $\theta_{j+1}\in\initball'$,
$\frac\stepsize2\sum_{i\le j}\|\nabla\Lout(\theta_i)\|^2\le\Lout(\theta_0)-\Lout(\theta_{j+1})\le C_2\initscale^2\mixrate^n$. By the
triangle and Cauchy--Schwarz inequalities,
\[
 \|\theta_{j+1}-\theta_0\|\le\stepsize\sum_{i\le j}\|\nabla\Lout(\theta_i)\|
 \le\bigl(2(j+1)\stepsize\,C_2\initscale^2\mixrate^n\bigr)^{1/2}\le\ballradius\initscale
\]
by the choice of $c_{\mathrm{time}}$. This closes the induction, and \cref{eq:lo-floor} gives the loss bound.

The descent inequality also gives $\Lout(\theta_j)\le\Lout(\theta_0)=\log K$, so \cref{eq:lo-identity} yields
$\E\,\KL(\bm u\|\hat y_j)\le\outcredit(\theta_j)\le C_2\initscale^2\mixrate^n$. Since $\hat y_{j,y}-1/K\le\mathrm{TV}(\hat y_j,\bm u)$,
Pinsker's and Jensen's inequalities give
\[
 \E\,\hat y_{j,y}-\frac1K\le\E\sqrt{\tfrac12\KL(\bm u\|\hat y_j)}\le\sqrt{\tfrac12\E\,\KL(\bm u\|\hat y_j)}\le\sqrt{C_2/2}\;\initscale\mixrate^{n/2}.
\]
Finally $\pimin/4\le1/4<\log2-1/4\le\log K-C_2\initscale_0^2\le\Lout(\theta_j)$, so $\theta_j\notin\basin(\Lout)$. The theorem
holds with $c=c_{\mathrm{time}}$, step-size bound $\stepsize_0$ and $C=\max(C_2,\sqrt{C_2/2})$.
\end{proof}

\begin{remark}[A zero readout leaves \cref{thm:separation} unchanged]
\label{rem:zero-readout}
Both formats can therefore share the initialization of \cref{thm:outcome-end}.
Part~(i) of \cref{thm:separation} uses $\hat y-\bm u=O(\initscale^2)$, which holds exactly at $\theta_0$ when $\tilde V=0$.
Part~(ii) bounds $\rulecred$ from above on $\initball$ for every readout. Its lower bound uses $F_n$ of
\cref{lem:leading}, a derivative in one readout entry equal to $\E_t\E[(\bm g_t)_kh^2_{i_t,l}]$. The residual $h^2$ does
not involve $V$, and $P_n$ depends only on attention and embedding entries, so fixing $\tilde V=0$ changes neither
the law of $P_n$ nor the event in the lower bound.
\end{remark}



\begin{remark}[Gradient flow]
\label{rem:flow}
The bounds of \cref{thm:outcome-end} also hold for gradient flow $\dot\theta=-\nabla\Lout(\theta)$ from $\theta_0$, with
$j\stepsize$ replaced by the time $t$, for all $t\le c\,\mixrate^{-n}$. The proof is the same, with the descent lemma replaced
by $\frac{d}{dt}\Lout(\theta(t))=-\|\nabla\Lout(\theta(t))\|^2$: while $\theta$ stays in $\initball$,
$\int_0^t\|\nabla\Lout\|^2\,ds\le C\initscale^2\mixrate^n$ by \cref{eq:lo-floor}, and
$\|\theta(t)-\theta_0\|\le(t\int_0^t\|\nabla\Lout\|^2\,ds)^{1/2}$ by Cauchy--Schwarz, so the first exit time from
$\initball$ is at least $c\,\mixrate^{-n}$. This uses only that $\Lout$ is smooth on $\initball'$, which the fixed ReLU
pattern gives.
\end{remark}

\begin{remark}[Width: an open point]
\label{rem:end-width}
(A1) fixes the widths, while the outcome executor needs MLP width $3K(n+1)$ and $d=2n+6$ positions.
\Cref{thm:outcome-end} holds at every fixed width. To conclude that an exact outcome solution exists in the same
class but is not reached, the constants $C_1$, $C_2$ and $\attconst$ must be tracked in the widths, through $a_0$,
the head dimension and the coefficient counts of \cref{lem:softmax-all,lem:size}. Polynomial growth in the width
would suffice, since the time bound is exponential in $n$. We have not carried out this tracking.
\end{remark}


\subsection{Technical verification of the score-count bounds}
\label{app:network}
\label{app:two-block}
\label{app:envelope}

This subsection proves the model statement in \cref{app:model-spectrum}.
We use normalized derivatives, Taylor coefficients and input support.

\paragraph{Setting.} Assume (A1) and (A2). Constants depend on the widths, $K$ and $\|\tilde\theta\|$, never
on $n$. Fix $a_0\ge1$ such that on $\initball$ every parameter block has operator norm at most $a_0\initscale$,
and every input $x_j=E\bm e_{{\mathrm{tok}}_j}+P\bm e_j$ has $\|x_j\|\le a_0\initscale$. A \emph{site} is a position whose
token varies: the letters in \cref{thm:main}, and in \cref{thm:separation} also $s_0$, the query
and the displayed states. For a parameter direction $\bm v$, $\partial_{\bm v}$ denotes the directional
derivative; a bound $\|\partial_{\bm v}\bm z\|\le a\|\bm v\|$ for all $\bm v$ gives $\|\Jz\|_F\le\sqrt K\,a$.

\paragraph{Affine MLPs.} A pre-activation is $b_m+(W_{\mathrm{in}}h)_m$.
On $\initball$, $|b_m|\ge\initscale/2$ and $|(W_{\mathrm{in}}h)_m|\le C\initscale^2$.
These bounds hold under score scaling for every $\xi\in[0,1]$.
For $\initscale<\initscale_0:=1/(2C)$, the sign is fixed by $\tilde b_m$.
In each block, a fixed diagonal $0/1$ matrix $D$ describes the ReLUs.
The network is smooth in $(\xi,\theta)$.
Block $\ell$ becomes

\begin{gather*}
 h^\ell_i=M_\ell\bigl(h^{\ell-1}_i+\bm a^\ell_i\bigr)+\bm o_\ell,
 \qquad
 \bm a^\ell_i=\sum_{\mathrm{heads}}\sum_{j\le i}\alpha_{ij}W_Vh^{\ell-1}_j,
 \\
 M_\ell=I+W^\ell_{\mathrm{out}}DW^\ell_{\mathrm{in}},\quad \bm o_\ell=W^\ell_{\mathrm{out}}D\bm b^\ell,
\end{gather*}
with $h^0_i=x_i$, $\bm z_i=Vh^2_i$, $\alpha_{i\cdot}=\mathrm{softmax}(\xi\psi_{i\cdot})$ and $\psi_{il}=(W_Qh_i)^\top W_Kh_l/\sqrt{d_{\mathrm{h}}}$.


\paragraph{Normalized derivatives.}
For $\ell=1,2$ and $k\ge0$, define
\[
 H_{\ell,k}=\sup_{\xi,i}\frac{\|\partial_\xi^kh_i^\ell(\xi)\|}{k!},\quad
 \Psi_{\ell,k}=\sup_{\xi,i,j}\frac{|\partial_\xi^k\psi_{ij}^\ell(\xi)|}{k!},\quad
 Z_k=\sup_{\xi,i}\frac{\|\partial_\xi^k\bm z_i(\xi)\|}{k!}.
\]
A dot adds $\partial_{\bm v}$ and the supremum over $\|\bm v\|=1$.
All suprema also cover $\initball$ and the inputs.

\begin{lemma}[Uniform softmax derivative bound]
\label{lem:softmax-all}
For $\alpha=\softmax(\psi)$ on any number of positions,
\begin{equation}
 \|\partial^k\alpha(\psi)[v_1,\ldots,v_k]\|_1
 \le 2^k k!\prod_{a=1}^k\|v_a\|_\infty,\qquad k\ge1.
 \label{eq:softmax-avg}
\end{equation}
The same bound times $\max_j\|u_j\|$ holds for
$\|\sum_j \partial^k\alpha_j[v_1,\ldots,v_k]u_j\|$.
\end{lemma}

\begin{proof}
Write $\langle f\rangle_\alpha=\sum_j\alpha_jf_j$.
\[
 \partial\alpha(\psi)[v]_j=\alpha_j(v_j-\langle v\rangle_\alpha),\qquad
 \partial_v\langle f\rangle_\alpha
 =\langle vf\rangle_\alpha-\langle v\rangle_\alpha\langle f\rangle_\alpha.
\]
After $k$ derivatives, each summand is a product of at most $k+1$ averages.
Each $v_a$ and $f$ occurs once. If $a_k$ is the sum of absolute coefficients, then
\[
 a_0=1,\qquad a_{k+1}\le2(k+1)a_k,\qquad a_k\le2^k k!.
\]
Thus
\[
 |\partial_{v_1}\cdots\partial_{v_k}\langle f\rangle_\alpha|
 \le2^k k!\|f\|_\infty\prod_a\|v_a\|_\infty.
\]
Take the supremum over $\|f\|_\infty\le1$ for the $\ell_1$ bound.
For vector-valued $f$, one average in each product is vector-valued; the same triangle inequality applies.
\end{proof}

\begin{lemma}[Normalized product and composition bounds]
\label{lem:classes}
Write $f_k=f^{(k)}/k!$. For a bilinear product of norm at most $1$,
\[
 \|(fg)_k\|\le\sum_{a+b=k}\|f_a\|\|h_b\|.
\]
If $\|\partial^jF\|/j!\le a_0 c^j$ and $\|h_k\|\le a b^k$ for $k\ge1$, then
\[
 \|(F\circ h)_k\|\le a_0ca\,b^k(1+ca)^{k-1},\qquad k\ge1.
\]
\end{lemma}

\begin{proof}
Leibniz gives the first bound. The chain rule in normalized form gives
\begin{align*}
 (F\circ h)_k
 &=\sum_{j=1}^k\frac1{j!}\partial^jF(h)
       \sum_{\substack{k_1+\cdots+k_j=k\\k_a\ge1}}
       [h_{k_1},\ldots,h_{k_j}],\\
 \|(F\circ h)_k\|
 &\le a_0 b^k\sum_{j=1}^k\binom{k-1}{j-1}(ca)^j
 =a_0ca\,b^k(1+ca)^{k-1}.
\end{align*}
\end{proof}

\paragraph{Absolute coefficient sums.}
For a fixed expansion $F=\sum_\nu w_\nu\term_\nu$ at $\xi=0$, let
\[
 \|F\|_{\mathrm{abs}}=\sum_\nu|w_\nu|N(\term_\nu),\qquad
 \|\partial F\|_{\mathrm{abs}}=\sum_\nu|w_\nu|
       \sup_{\initball,\,\|\bm v\|=1}\|\partial_{\bm v}\term_\nu\|.
\]
Here $N(\term)$ multiplies the $a_0\initscale$ bounds on parameter and input factors.
Both sums bound the norm before cancellation. The first is subadditive and submultiplicative;
 the tangent sum obeys the product rule.
They are needed when summing over input supports.


\begin{lemma}[Score-count size bound]
\label{lem:size}
There are $\initscale_0>0$, $C$ and $C_{\mathrm{size}}$, independent of $n$ and $k$, such that for $\initscale\le\initscale_0$, $\theta\in \initball$, every position $i$,
every input and every $k\ge1$:
\begin{enumerate}[label=(\alph*),leftmargin=*,itemsep=0pt,topsep=2pt]
 \item $\|h^\ell_i-x_i\|\le C\initscale^2$, $|\psi_{il}|\le C\initscale^4$, $\|\bm z_i\|\le C\initscale^2$ and $\|\partial_{\bm v}\bm z_i\|\le C\initscale\|\bm v\|$, at every $\xi\in[0,1]$.
 \item For every $\xi\in[0,1]$,
 \[
 \frac{\|\bm z_i^{(k)}(\xi)\|}{k!}\le C\initscale^3(C_{\mathrm{size}}\initscale^4)^k,
 \qquad
 \frac{\|\partial_{\bm v}\bm z_i^{(k)}(\xi)\|}{k!}\le C\initscale^2(C_{\mathrm{size}}\initscale^4)^k\|\bm v\|,
 \]
 and the same for $h^2_i$ with $C\initscale^2,C\initscale$ in place of $C\initscale^3,C\initscale^2$. Hence $\|\bm z^{{\mathrm{rem}},\bar k}_i\|\le C\initscale^3(C_{\mathrm{size}}\initscale^4)^{\bar k+1}$ and
 $\|\partial_{\bm v}\bm z^{{\mathrm{rem}},\bar k}_i\|\le C\initscale^2(C_{\mathrm{size}}\initscale^4)^{\bar k+1}\|\bm v\|$ for every $\bar k\ge0$.
 \item $\|\bm z_i^{[k]}\|_{\mathrm{abs}}\le C\initscale^3(C_{\mathrm{size}}\initscale^4)^k$ and $\|\partial\bm z_i^{[k]}\|_{\mathrm{abs}}\le C\initscale^2(C_{\mathrm{size}}\initscale^4)^k$,
 and the same holds for $h^{2,[k]}_i$ with $C\initscale^2,C\initscale$ in place of $C\initscale^3,C\initscale^2$.
\end{enumerate}
\end{lemma}


\begin{proof}
\proofstep{(a) Amplitudes.}
For $\ell=1,2$,
\begin{align*}
 \|\bm a_i^\ell\|&\le3a_0\initscale\max_j\|h_j^{\ell-1}\|\le C\initscale^2,\\
 \|M_\ell-I\|+\|\bm o_\ell\|&\le C\initscale^2,\qquad
 H_{\ell,0}\le C\initscale,\qquad\|h_i^\ell-x_i\|\le C\initscale^2,\\
 \Psi_{\ell,0}&\le C\initscale^2H_{\ell-1,0}^2\le C\initscale^4,
 \qquad Z_0\le\|V\|H_{2,0}\le C\initscale^2.
\end{align*}
Unit parameter directions satisfy
\[
 \dot H_{0,0}\le C,\quad
 \dot \Psi_{\ell,0}\le C\initscale^3,
 \quad \|\partial_{\bm v}\alpha_{i\cdot}^\ell\|_1\le C\initscale^3,
 \quad \dot H_{\ell,0}\le C,\quad \dot Z_0\le C\initscale.
\]

\proofstep{(b) Derivatives and remainder.}
Let $\Xi_{\ell,k}$ be the normalized derivative bound for the scaled score $\xi\psi^\ell$.
Leibniz and \cref{lem:softmax-all,lem:classes} give
\begin{align*}
 \Psi_{\ell,k}&\le C\initscale^2\sum_{a+b=k}H_{\ell-1,a}H_{\ell-1,b},\\
 \Xi_{\ell,k}&\le \Psi_{\ell,k}+\Psi_{\ell,k-1}\quad(k\ge1),\\
 U_{\ell,k}&:=\sup_{\xi,i}\frac{\|\partial_\xi^k\alpha_{i\cdot}^\ell\|_1}{k!}
 \le\sum_{j=1}^k2^j
    \sum_{\substack{a_1+\cdots+a_j=k\\a_r\ge1}}
    \prod_{r=1}^j\Xi_{\ell,a_r},\\
 B_{\ell,k}&:=\sup_{\xi,i}\frac{\|\partial_\xi^k\bm a_i^\ell\|}{k!}
 \le C\initscale\sum_{a+c=k}U_{\ell,a}H_{\ell-1,c},\\
 H_{\ell,k}&\le C(H_{\ell-1,k}+B_{\ell,k}),\qquad Z_k\le C\initscale H_{2,k}.
\end{align*}
Here $U_{\ell,0}=1$ and $H_{0,k}=0$ for $k\ge1$.
For a fixed $b=C\initscale^4$, the first block gives
\[
 U_{1,k}\le b^k,\qquad
 B_{1,k}+H_{1,k}\le C\initscale^2b^k\quad(k\ge1).
\]
The second block then satisfies, after enlarging fixed factors,
\begin{align*}
 \Psi_{2,k}&\le C\initscale^5(k+1)b^k\le C\initscale^5(2b)^k,\\
 \Xi_{2,k}&\le C(2b)^k,\\
 U_{2,k}&\le C\sum_{j=1}^k\binom{k-1}{j-1}C^j(2b)^k
                \le C[C b]^k,\\
 B_{2,k}+H_{2,k}&\le C\initscale^2[C b]^k,\qquad
 Z_k\le C\initscale^3(C_{\mathrm{size}}\initscale^4)^k.
\end{align*}
Each $C$ here is fixed; only two blocks are composed.

For tangents, differentiate the preceding recurrences. In particular,
\begin{align*}
 \dot \Psi_{\ell,k}
 &\le C\initscale\sum_{a+c=k}H_{\ell-1,a}H_{\ell-1,c}
     +C\initscale^2\sum_{a+c=k}
       (\dot H_{\ell-1,a}H_{\ell-1,c}+H_{\ell-1,a}\dot H_{\ell-1,c}),\\
 \dot\Xi_{\ell,k}&\le\dot \Psi_{\ell,k}+\dot \Psi_{\ell,k-1},\\
 \dot B_{\ell,k}
 &\le C\sum_{a+c=k}U_{\ell,a}H_{\ell-1,c}
      +C\initscale\sum_{a+c=k}
       (\dot U_{\ell,a}H_{\ell-1,c}+U_{\ell,a}\dot H_{\ell-1,c}).
\end{align*}
A differentiated softmax term has either one $\dot\Xi$ factor or
$\partial^{j+1}\alpha[\dot\Xi_{\ell,0},\ldots]$.
By \cref{lem:softmax-all}, the extra factor $j+1\le2^k$ is absorbed into the derivative scale.
The two-block induction yields
\[
 \dot U_{\ell,k}\le C\initscale^{-1}(C_{\mathrm{size}}\initscale^4)^k,\quad
 \dot B_{\ell,k}+\dot H_{\ell,k}\le C\initscale(C_{\mathrm{size}}\initscale^4)^k,\quad
 \dot Z_k\le C\initscale^2(C_{\mathrm{size}}\initscale^4)^k.
\]
Taylor's integral remainder gives
\begin{align*}
 \|\bm z_i^{{\mathrm{rem}},\bar k}\|
 &\le\sup_\xi\frac{\|\partial_\xi^{\bar k+1}\bm z_i(\xi)\|}{(\bar k+1)!}
 \le C\initscale^3(C_{\mathrm{size}}\initscale^4)^{\bar k+1},\\
 \|\partial\bm z_i^{{\mathrm{rem}},\bar k}\|
 &\le C\initscale^2(C_{\mathrm{size}}\initscale^4)^{\bar k+1}.
\end{align*}
The same estimates before the readout give the residual bounds.

\proofstep{(c) Coefficient sums.}
At $\xi=0$, every average in the softmax proof is uniform, and
\[
 \left\|\frac1m\sum_j\prod_rF_r(j)\right\|_{\mathrm{abs}}
 \le\prod_r\max_j\|F_r(j)\|_{\mathrm{abs}}.
\]
Thus the recurrences in (b) also bound $\|\bm z_i^{[k]}\|_{\mathrm{abs}}$.
For a monomial with $f(\term)$ parameter and input occurrences,
\[
 \sup\|\partial_{\bm v}\term\|
 \le\frac{2f(\term)}{a_0\initscale}N(\term)\|\bm v\|,
 \qquad f(\term)\le C(k+1),\qquad k+1\le2^k.
\]
The occurrence bound follows by induction over the two blocks and the score factors.
Enlarging $C_{\mathrm{size}}$ gives $\|\partial\bm z_i^{[k]}\|_{\mathrm{abs}}\le C\initscale^2(C_{\mathrm{size}}\initscale^4)^k$.
\end{proof}

\begin{lemma}[Score-count support bound]
\label{lem:support}
For every position $i$ and every $k\ge1$:
\begin{enumerate}[label=(\alph*),leftmargin=*,itemsep=0pt,topsep=2pt]
 \item $\bm z^{[0]}_i=\bm z(0)$ is affine in the inputs: a sum of terms that each contain one input $x_j$
 or none, with total weight at most $C$. A transported input term has $\|\partial_{\bm v}\term\|\le C\initscale^2\|\bm v\|$;
 the direct input term and input-free terms have bound $C\initscale\|\bm v\|$.
 \item Every term of $\bm z^{[k]}_i$, and of $h^{\ell,[k]}_i$, contains at most $2k+1$ inputs.
 \item The one-score part is the first-order term of the softmax, displayed below; each of its terms $\term$
 has $\|\term\|\le C\initscale^7$ and $\|\partial_{\bm v}\term\|\le C\initscale^6\|\bm v\|$.
\end{enumerate}
\end{lemma}


\begin{proof}
\proofstep{(a) Zero scores.}
With $W_V^\ell$ denoting the sum of value maps,
\begin{align*}
 h_j^1(0)&=M_1\left(x_j+\frac1{m_j}\sum_{l\le j}W_V^1x_l\right)+\bm o_1,
 \\
 \bm z_i^{[0]}&=V\left[M_2\left(h_i^1(0)+\frac1{m_i}\sum_{j\le i}W_V^2h_j^1(0)\right)+\bm o_2\right].
\end{align*}
Every term is affine in the inputs. A transported term contains $V,W_V,x_j$;
its tangent is $O(\initscale^2)$. The direct term $VM_2M_1x_i$ has tangent $O(\initscale)$.

\proofstep{(b) Input support.}
Let $d(F)$ count the largest number of input occurrences in an expanded monomial.
\[
 d(F+G)\le\max(d(F),d(G)),\qquad d(FG)\le d(F)+d(G),\qquad d(x_i)=1.
\]
Starting from $d(h^{0,[0]})=1$, induction gives
\begin{align*}
 d(h^{\ell,[a]})&\le2a+1,\\
 d(\psi^{\ell,[b]})
 &\le\max_{a+c=b}\{(2a+1)+(2c+1)\}=2b+2,\\
 d((\xi\psi^\ell)^{[a]})&\le2a\qquad(a\ge1).
\end{align*}
A head coefficient of order $k$ has $a_1+\cdots+a_j+c=k$, so
\[
 d\left(\frac1{j!}\partial^j\alpha(0)
       [\psi^{[a_1-1]},\ldots,\psi^{[a_j-1]}]W_Vh^{[c]}\right)
 \le2a_1+\cdots+2a_j+2c+1=2k+1.
\]
Affine MLPs and the readout introduce no new input products.
Parameter differentiation cannot increase this count.

\proofstep{(c) One score.}
The contribution from differentiating one attention weight is
\[
 \frac1{m_i}\sum_{j\le i}(\psi_{ij}(0)-\bar\psi_i(0))W_Vh_j(0),
 \qquad \bar\psi_i(0)=\frac1{m_i}\sum_{l\le i}\psi_{il}(0).
\]
Each logit term contains $V,W_V,W_Q,W_K$ and three residual inputs. Hence
\[
 \|\term\|\le C\initscale^7,\qquad
 \|\partial_{\bm v}\term\|\le C\initscale^6\|\bm v\|.
\]
\end{proof}


\paragraph{Proof of \cref{prop:score-degree}.}
\begin{proof}
For a term with site set $U$, the Efron--Stein formula gives
\[
 \term_S=0\quad(S\nsubseteq U),\qquad
 \sup\|\term_S\|\le2^{|S|}\sup\|\term\|,
 \qquad \sum_{S\subseteq U}2^{|S|}=3^{|U|}.
\]
At the answer position, the direct input is fixed. By \cref{lem:support}(a),
\[
 \sup\|J^{[0]}\|_F\le C\initscale,\qquad
 \sum_{|S|=1}\sup\|J_S^{[0]}\|_F\le C\initscale^2,\qquad
 J_S^{[0]}=0\quad(|S|\ge2).
\]
For $k\ge1$, \cref{lem:size,lem:support} give
\[
 \sum_S\sup\|J_S^{[k]}\|_F
 \le\sqrt K\,3^{2k+1}\|\partial\bm z^{[k]}\|_{\mathrm{abs}}
 \le C\initscale^2(\attconst\initscale^4)^k,\qquad
 J_S^{[k]}=0\quad(|S|>2k+1),
\]
with $\attconst\ge9C_{\mathrm{size}}$. The remainder bound is \cref{lem:size}(b).
The logit bounds have one additional factor $\initscale$.
\end{proof}


\begin{lemma}[Position-selection penalty]
\label{lem:dilution}
Fix positions $l<i$ and $k\ge1$, and let
\begin{equation}
 \kappa_{l,i}:=\frac1{m_i}\Bigl(1+\sum_{l<j\le i}\frac1{m_j}\Bigr).
 \label{eq:dilution}
\end{equation}
After enlarging $\attconst$, the terms of $\bm z^{[k]}_i$ that contain $x_l$ have $\|\cdot\|_{\mathrm{abs}}$ at most $C\initscale^3(\attconst\initscale^4)^k\kappa_{l,i}$ and $\|\partial(\cdot)\|_{\mathrm{abs}}$ at most $C\initscale^2(\attconst\initscale^4)^k\kappa_{l,i}$, and the same holds for $h^{2,[k]}_i$ with $C\initscale^2,C\initscale$ in place of
$C\initscale^3,C\initscale^2$.
\end{lemma}


\begin{proof}
Let $[F]_l$ collect expanded terms containing $x_l$.
Uniform attention in block 1 gives
\[
 \frac{\|[h_j^{1,[0]}]_l\|_{\mathrm{abs}}}{C\initscale}
 \le\begin{cases}0,&j<l,\\1,&j=l,\\C\initscale/m_j,&j>l.\end{cases}
\]
The same upper bounds hold for the relative tangent sums.
In block 2, replacing one residual factor by its $x_l$ part costs at most
\[
 \frac{C\initscale}{m_i}\quad\text{(query residual)},\qquad
 \frac1{m_i}\quad\text{(key $l$)},\qquad
 \frac{C\initscale}{m_i}\sum_{l<j\le i}\frac1{m_j}\quad\text{(later keys)}.
\]
The residual connection contributes at cost $C/m_i$.
There are at most $C(k+1)$ residual-factor choices in an order-$k$ term. Therefore
\begin{align*}
 \|[\bm z_i^{[k]}]_l\|_{\mathrm{abs}}
 &\le C(k+1)\initscale^3(C_{\mathrm{size}}\initscale^4)^k\kappa_{l,i}
 \le C\initscale^3(\attconst\initscale^4)^k\kappa_{l,i},\\
 \|\partial[\bm z_i^{[k]}]_l\|_{\mathrm{abs}}
 &\le C\initscale^2(\attconst\initscale^4)^k\kappa_{l,i}.
\end{align*}
The same induction before the readout gives the residual estimates.
\end{proof}

%=============================================================================


\section{Numerical verification of the theory}
\label{app:verify}

\Cref{tab:verify} maps claims to exact computations.
Gradients and spectra are summed over all words in float64.
Unless stated otherwise, the bias condition (A2) holds.
The tables below give detailed measurements.

\begin{table}[htbp]
\centering\small
\fitwidth{\begin{tabular}{@{}rllll@{}}
\toprule
& Claim & What is measured & Predicted & Measured\\
\midrule
1 & \Cref{eq:horizon}, counting & answer-dependent gradient vs $n$ & factor $\mixrate=0.5$ per letter & $0.48$--$0.51$ (\cref{tab:spectrum})\\
2 & \Cref{eq:P1}, degree $r$ & exponent of $\sum_{|S|=r}\|(\Jz)_S\|$ in $\initscale$ & $2$; $6$ ($r=2,3$); $\ge10$ & $2.0,6.1,7.0,11.1,12.0$ for $r=1..5$\\
3 & \Cref{eq:P1} gives $\sum_{r\ge2}J_r\le C\initscale^{12}$ & exponent of $\Jhigh/\Jtotal$ in $\initscale$ & $\ge10$ & $10.00$--$10.11$ ($n=4,5$)\\
4 & \Cref{prop:plateau} & $\|\pd{\Lout}\theta(\theta_0)\|/\initscale(\initscale^2+\mixrate^n+\initscale\mixrate^{n-1})$, $n=4..10$ & bounded & $3.9$--$5.3$\\
5 & \Cref{thm:separation}(i) & $\|\pd{\Lstep^{\ge2}}V\|/\initscale$, $n=4,6,8$ & $\|\bm\Gamma\tilde E_{\mathrm{val}}\|_F=0.811$ & $0.807$--$0.816$\\
6 & \Cref{thm:separation}(iii), \cref{lem:leading} & rule credit at $\tilde\theta^\circ$ vs $\initscale^6P_n(\tilde\theta^\circ)$ & equal to leading order & equal to six digits, $n=2..12$\\
7 & \Cref{thm:separation}(ii)--(iii) & $n\rulecred/\initscale^6$, fixed query & flat in $n$ & \cref{tab:verify-query}\\
8 & \Cref{prop:represent} & executors rebuilt from \cref{app:params} & identical tensors, exact & identical; exact up to $n=1000$\\
\bottomrule
\end{tabular}}
\caption{\textbf{Credit and plateau bounds are checked against exact computations.}
All rows use the two-block Transformer of (A1) at width $16$, counting with $K=2$ unless
stated; exponents in $\initscale$ are fitted on $[0.005,0.05]$.}
\label{tab:verify}
\end{table}

\paragraph{Row 2.} Two blocks, width $16$, $n=6$, exact over all $4^6$ words. The proved exponents are
attained at $r=1,2$, the degrees the theorems use.

\paragraph{Rows 4 and 5.} $K=2$, width $16$, exact over all words and both start states. Over $20$ updates at
$\stepsize=0.05$, for $n=4,6,8$ and $\initscale\in\{10^{-2},3\times10^{-3}\}$, the step loss falls by $1.5$--$2.2\,\initscale^2$,
independent of $\initscale$ and $n$, while the outcome loss falls by $7\times10^{-9}$--$7\times10^{-6}\,\initscale^2$.

\begin{table}[htbp]
\centering\small
\begin{tabular}{@{}lccccc@{}}
\toprule
$n$ & $4$ & $5$ & $6$ & $7$ & $8$\\
\midrule
fixed query, $n\rulecred/\initscale^6$, draws 1, 2 & $309$, $133$ & $308$, $130$ & $309$, $130$ & $324$, $131$ & $304$, $134$\\
\bottomrule
\end{tabular}
\caption{\textbf{Row 7: the rule credit is $\Theta(\initscale^6/n)$ for a fixed query.}
$|\Sigma|=4$ uniform letters, $p=\frac14$, $\mixrate=\frac12$, $\initscale=0.02$, exponent in $\initscale$ $6.00$ throughout.}
\label{tab:verify-query}
\end{table}


%=============================================================================
\section{Experiments}
\label{app:experiments}
\label{app:experimental-details}

The subsections follow the predictions in \cref{sec:validation}.

\subsection{Protocol}
\label{app:letter-protocol}

\begin{center}\small
\fitwidth{\begin{tabular}{@{}llll@{}}
\toprule
& Letter tasks (two-block) & Boolean-8, register (GPT-2L) & LoRA (SmolLM2, Qwen2.5)\\
\midrule
Model & \cref{app:model}, $\blocks=2$, $3$ heads & $2$ blocks, $4$ heads, width $128$, LayerNorm & rank $8$\\
Data & $20{,}000$ train, $1{,}000$ test & $20{,}000$ train, $500$--$1{,}000$ test & $12{,}000$ train\\
Initialization & every tensor $\mathrm N(0,0.02^2)$ & default & default LoRA\\
Optimizer & AdamW, lr $10^{-3}$, clip $1$, batch $256$ & AdamW, lr $3\times10^{-4}$, batch $128$ & batch $16$\\
Budget & $8{,}000$ updates ($16$k or $32$k where stated) & $8{,}000$ updates & $800$ updates\\
Seeds & $42$--$44$ ($42$--$51$ at $n=16$) & three or five & $2001$--$2003$\\
\bottomrule
\end{tabular}}
\end{center}
Paired formats share initialization, prompts, batches, optimizer and budget.
Only the supervised continuation differs.
We score greedy decoding on held-out prompts.
Letter tasks use four letters; first-occurrence words contain the query.
For GPT-2L and LoRA, seeds change only initialization.
Their spread does not capture data variability.

\subsection{Experiments inside the theorems' regime}
\label{app:exp-regime}
\label{app:regime}

These runs satisfy (A1)--(A4) and give \cref{fig:regime,tab:matched}. The runs of \cref{fig:regime} and the first
table start from a zero readout, as in \cref{thm:outcome-end}; the other measurements here keep the nonzero readout of
\cref{prop:plateau}.
Use $K=2$, two blocks, width $16$, three heads of dimension $8$, and MLP width $16$.
Every initial scaled bias has $|\tilde b_m|\ge1$.
Compute full-batch gradients in float64; the step objective is $\Lstep^{\ge2}$.
Two letters with match probability $1/4$ allow enumeration up to $n=14$.
They give the same $\mixrate=1/2$ as four uniform letters.
Each update moves the parameters by $10^{-3}\initscale\|\tilde\theta\|$ along $-\nabla L$, and the training time is the sum
of the step sizes. Define $\traintime(\ballradius)$ by $\|\theta_j-\theta_0\|=\ballradius\initscale\|\tilde\theta\|$.
These runs approximate gradient flow (\cref{rem:flow}): at $n=12$, halving or quartering the update length changes
$\traintime(0.1)$ by less than $1\%$ for both losses.
No ReLU switches with the letters in any run.

\paragraph{The outcome loss takes time of order $\mixrate^{-n}$ to leave the ball.} From a zero readout, the outcome time
grows by $4.4$ to $5.1$ per two letters at $\initscale=0.002$, against $\mixrate^{-2}=4$, as \cref{thm:outcome-end} predicts.
At $\initscale=0.02$ it is shorter at small $n$ and meets the $\initscale=0.002$ times by $n=14$. The step time does not
depend on $n$; in that time the step loss falls by $10$ to $13\initscale^2$ at every $n$, which is progress on the prior,
not yet the rule. With the nonzero readout of \cref{prop:plateau}, the outcome times are within a factor of $2$ of these,
so they also grow past its cap $\initscale^{-2}$: the answer-independent part of the gradient removes the bias of the
logits and then vanishes.
\begin{center}\small
\begin{tabular}{@{}lcccccc@{}}
\toprule
$\traintime(0.1)$, zero readout & $n=4$ & $6$ & $8$ & $10$ & $12$ & $14$\\
\midrule
outcome, $\initscale=0.002$, draw 1 & $9.3\mathrm{e}3$ & $4.6\mathrm{e}4$ & $2.3\mathrm{e}5$ & $1.0\mathrm{e}6$ & $4.7\mathrm{e}6$ & $2.2\mathrm{e}7$\\
outcome, $\initscale=0.002$, draw 2 & $1.2\mathrm{e}4$ & $6.1\mathrm{e}4$ & $2.9\mathrm{e}5$ & $1.3\mathrm{e}6$ & $6.2\mathrm{e}6$ & $2.8\mathrm{e}7$\\
outcome, $\initscale=0.002$, draw 3 & $1.5\mathrm{e}4$ & $7.7\mathrm{e}4$ & $3.8\mathrm{e}5$ & $1.8\mathrm{e}6$ & $7.8\mathrm{e}6$ & $3.8\mathrm{e}7$\\
outcome, $\initscale=0.02$, draw 1 & $7.8\mathrm{e}2$ & $3.8\mathrm{e}3$ & $1.9\mathrm{e}4$ & $9.4\mathrm{e}4$ & $2.1\mathrm{e}6$ & $3.0\mathrm{e}7$\\
step, $\initscale=0.002$, draws 1, 2, 3 & \multicolumn{6}{c}{$4.18$, $4.33$, $5.55$ at every $n$ (to three digits)}\\
step, $\initscale=0.02$, draw 1 & $4.17$ & $4.18$ & $4.19$ & $4.19$ & $4.20$ & $4.20$\\
\bottomrule
\end{tabular}
\end{center}
With four letters and $\initscale=0.02$, the initial outcome gradient is $3.9$ to $5.3$ times $\initscale(\initscale^2+\mixrate^n+\initscale\mixrate^{n-1})$ for
$n=4,\ldots,10$, and $(\log K-\Lout)/\initscale^2$ at $\traintime(0.1)$ falls from $2.46$ to $0.010$, as \cref{prop:plateau}
predicts.

\paragraph{The exit time depends on $\mixrate$ and $n$ through $\initscale^2+\mixrate^n$.} Varying the match probability
$p\in\{0.1,0.15,0.25,0.35\}$ separates $\mixrate=1-2p$ from $n$ ($\initscale=0.002$, first draw). The table gives
$\traintime(0.1)(\initscale^2+\mixrate^n+\initscale\mixrate^{n-1})$, which \cref{prop:plateau} bounds below by a constant, since $\initscale\le\mixrate$. Across $\mixrate\in\{0.5,0.7,0.8\}$ it
varies by at most $3\%$ at each $n$, while $\traintime(0.1)$ itself varies by up to a factor of $285$. At $\mixrate=0.3$
and $n\ge8$, $\mixrate^n$ falls to the order of $\initscale^2$ and the time exceeds the bound.
\begin{center}\small
\begin{tabular}{@{}lccccc@{}}
\toprule
& $n=4$ & $n=6$ & $n=8$ & $n=10$ & $n=12$\\
\midrule
$\mixrate=0.8$ & $482$ & $602$ & $724$ & $841$ & $972$\\
$\mixrate=0.7$ & $482$ & $603$ & $726$ & $841$ & $981$\\
$\mixrate=0.5$ & $483$ & $607$ & $731$ & $845$ & $1{,}002$\\
$\mixrate=0.3$ & $485$ & $611$ & $790$ & $1{,}741$ & $26{,}093$\\
\bottomrule
\end{tabular}
\end{center}

\paragraph{Curvature.} The top Hessian eigenvalue of $\outcredit$ at $\theta_0$ falls from $2.2\times10^{-3}$ at
$n=4$ to $1.0\times10^{-6}$ at $n=14$, while that of the step loss stays at $0.25$--$0.26$ ($\initscale=0.02$, Lanczos).

\paragraph{The rule credit is of order $\initscale^6/n$ and has no factor $\mixrate$.} Credits at initialization,
three draws of $\tilde\theta$, $\initscale=0.02$. Between $\initscale=0.005$ and $0.02$ the exponents in $\initscale$ are $1.00$ for the prior,
$6.00$ for the rule and $2.1$ for the total outcome credit, in every row. The total outcome credit is
the larger up to $n=14$, but almost all of it is at degree $0$, which no executor is built from.
\begin{center}\small
\fitwidth{\begin{tabular}{@{}lcccccc@{}}
\toprule
& $n=4$ & $6$ & $8$ & $10$ & $12$ & $14$\\
\midrule
prior credit, draw 1 & $3.3\mathrm{e}{-2}$ & $3.3\mathrm{e}{-2}$ & $3.3\mathrm{e}{-2}$ & $3.3\mathrm{e}{-2}$ & $3.3\mathrm{e}{-2}$ & $3.3\mathrm{e}{-2}$\\
rule credit, draw 1 & $3.3\mathrm{e}{-9}$ & $2.5\mathrm{e}{-9}$ & $1.8\mathrm{e}{-9}$ & $1.3\mathrm{e}{-9}$ & $1.1\mathrm{e}{-9}$ & $8.8\mathrm{e}{-10}$\\
rule credit $\times\,n/\initscale^6$, draws 1, 2, 3 & $207$, $195$, $120$ & $231$, $194$, $112$ & $227$, $175$, $97$ & $208$, $144$, $122$ & $211$, $141$, $145$ & $192$, $147$, $169$\\
total outcome credit, draw 1 & $1.6\mathrm{e}{-4}$ & $3.2\mathrm{e}{-5}$ & $6.6\mathrm{e}{-6}$ & $1.4\mathrm{e}{-6}$ & $3.2\mathrm{e}{-7}$ & $7.1\mathrm{e}{-8}$\\
\bottomrule
\end{tabular}}
\end{center}

\paragraph{The rule credit exceeds the outcome interaction credit by a margin that grows
exponentially.} We split the answer exactly into its Efron--Stein components over the letters,
for each start state, and measure the outcome credit at degrees $\ge2$ under a uniform start. The
interaction credit has exponent $7.0$ in $\initscale$ in every row (observed, not proved). With the start state
fixed, as in \cref{thm:main}, its exponent is $6.0$, and the rule credit exceeds it by $1.3$ to $2.5$
times at $n=4$ and $30$ to $70$ times at $n=12$.
\begin{center}\small
\fitwidth{\begin{tabular}{@{}lcccccc@{}}
\toprule
& $n=4$ & $6$ & $8$ & $10$ & $12$ & $14$\\
\midrule
interaction credit, uniform start, draw 1 & $2.8\mathrm{e}{-11}$ & $9.2\mathrm{e}{-12}$ & $2.6\mathrm{e}{-12}$ & $6.7\mathrm{e}{-13}$ & $1.7\mathrm{e}{-13}$ & $4.2\mathrm{e}{-14}$\\
rule$/$interaction, draws 1, 2, 3 & $119$, $64$, $80$ & $268$, $129$, $150$ & $709$, $315$, $351$ & $2{,}000$, $803$, $1{,}350$ & $6{,}730$, $2{,}600$, $5{,}310$ & $21{,}100$, $9{,}410$, $21{,}300$\\
\bottomrule
\end{tabular}}
\end{center}

\subsection{Predictions 1--3: spectra and gradients at initialization}
\label{app:exp-spectra}
\label{app:spectrum}
\label{app:validation-spectra}

\paragraph{Answer spectra.} For counting, $A_r$ has the closed form of \cref{rem:tight}. For GPT-2L tasks, propagate the second moment of $\bm v_t$
from \cref{lem:target-contraction}.
The resulting $\sum_r A_r$ matches the centred answer variance $1-1/K$.

Boolean-8 has $4$ bits ($K=16$) and NOT, CNOT, SWAP and Toffoli gates on uniformly
random wires ($52$ distinct gates, $n=8$). The register machine has $(x,y)\in\mathbb Z_{17}^2$ ($K=289$), five
instructions chosen uniformly, $n=16$, and answer $x$. Both mix slowly in the worst case, but almost
none of their answer is visible at degrees $0$ and $1$.

\begin{table}[htbp]
\centering\small
\fitwidth{\begin{tabular}{@{}lcccc@{}}
\toprule
Task & $\mixrate$ & share at $r\le1$ & share at $r\ge2$ & mean degree\\
\midrule
Count a letter mod $K$, four letters & $|\frac34+\frac14e^{2\pi i/K}|$ & $\le(1+n)\mixrate^{2(n-1)}$ & -- & linear in $n$\\
Boolean-8 ($n=8$) & $0.820$ & $1.1\%$ & $98.9\%$ & $6.3$\\
Register machine ($n=16$) & $0.953$ & $0.08\%$ & $99.9\%$ & $12.3$\\
First occurrence & -- & $\ge p(1-p)$ at degree $1$ & -- & --\\
\bottomrule
\end{tabular}}
\caption{\textbf{The tasks of the paper.} Share of the answer's variance by degree. First occurrence
is not a permutation task and keeps a degree-$1$ part at every $n$ (\cref{rem:control}).}
\label{tab:family}
\label{tab:validation-spectra}
\end{table}

\paragraph{Jacobian spectra.} \Cref{fig:overlap} uses the AdamW architecture and initialization.
Estimate $J_r$ from $\E\langle\Jz(x),\Jz(x_\varrho)\rangle_F=\sum_r\varrho^rJ_r$.
The coupling keeps each letter with probability $\varrho$ and otherwise redraws it.
At $n=8$, we also compute an exact transform:

for a Gaussian parameter direction $\bm g$, $\mathbb E_{\bm g}\|(\Jz\bm g)_r\|^2=\|(\Jz)_r\|_F^2$, and the
Efron--Stein split of $\Jz\bm g$ over the $4^8$ letter grid is an orthonormal transform along each
letter axis. We average $16$ directions over all eight pairs $(s_0,q)$. At the experiments' scale
$\initscale=0.02$, the share of Jacobian energy above degree $1$ is $2.4\times10^{-6}$ and $1.1\times10^{-5}$ for two seeds. It
grows steeply with $\initscale$ and exceeds $0.99$ at $\initscale=0.5$, so the low-degree Jacobian is a property of
small initialization, not of the architecture alone.

\paragraph{A pretrained model.} \Cref{tab:bbh-spectrum} repeats the exact split for Qwen2.5-0.5B with
rank-8 LoRA on three-object BIG-Bench Hard items, where the answer is pure degree $3$.

\begin{table}[htbp]
\centering\small
\begin{tabular}{@{}lcccc@{}}
\toprule
degree & $0$ & $1$ & $2$ & $3$\\
\midrule
answer, $A_r/\sum A$ & $0$ & $0$ & $0$ & $100\%$\\
LoRA gradient, $J_r/\Jtotal$ & $99.60\msd{}{0.02}\%$ & $0.256\msd{}{0.011}\%$ & $0.118\msd{}{0.007}\%$ & $0.022\msd{}{0.002}\%$\\
\bottomrule
\end{tabular}
\caption{\textbf{A pretrained model's gradient is also low-degree.} Each context fixes the scene,
items, start and query, and its three swaps range over all $6^3$ values. We take the full
Jacobian of the three option-letter logits with respect to the $2.4$M parameters of the LoRA
up-projection and split it exactly by degree. Mean over three LoRA initializations, ten contexts
each. The measured $\|\pd\outcredit\theta\|$ is $2$--$6\%$ of the bound of \cref{lem:gradient-split}.}
\label{tab:bbh-spectrum}
\end{table}

\paragraph{Prediction 3: the answer gradient decays at the predicted rate.} We sum gradients
exactly over all $8\cdot4^n$ prompts. A log-linear fit gives $0.508$ per letter ($R^2=0.995$) for seed $42$,
and seeds $43$ and $44$ give $0.513$ and $0.480$, against $\mixrate=0.5$.

\begin{table}[htbp]
\centering\small
\fitwidth{\begin{tabular}{@{}lcccccc@{}}
\toprule
$n$ & $2$ & $4$ & $6$ & $8$ & $9$ & $10$\\
\midrule
answer-dependent gradient $\|\mathbb E[(\Jz)^\top(y-\bm u)]\|$ & $8.6\mathrm{e}{-4}$ & $2.5\mathrm{e}{-4}$ & $6.5\mathrm{e}{-5}$ & $2.1\mathrm{e}{-5}$ & $7.6\mathrm{e}{-6}$ & $3.6\mathrm{e}{-6}$\\
step supervision gradient & $1.8\mathrm{e}{-2}$ & $3.0\mathrm{e}{-2}$ & $3.9\mathrm{e}{-2}$ & $4.0\mathrm{e}{-2}$ & $4.2\mathrm{e}{-2}$ & $4.3\mathrm{e}{-2}$\\
ratio & $21$ & $119$ & $598$ & $1{,}849$ & $5{,}561$ & $11{,}708$\\
full answer-token gradient & $0.159$ & $0.146$ & $0.165$ & $0.183$ & $0.143$ & $0.168$\\
$J_0$ & $0.111$ & $0.111$ & $0.148$ & $0.143$ & $0.120$ & $0.127$\\
$\sum_{r\ge1}J_r$ & $2.1\mathrm{e}{-4}$ & $3.0\mathrm{e}{-4}$ & $4.2\mathrm{e}{-4}$ & $4.6\mathrm{e}{-4}$ & $4.4\mathrm{e}{-4}$ & $4.0\mathrm{e}{-4}$\\
\bottomrule
\end{tabular}}
\caption{\textbf{The answer signal decays at the predicted rate, but the step signal does not.}
Two-block Transformer at initialization, counting with $K=2$ over four letters ($\mixrate=0.5$), seed
$42$. Almost all of the outcome gradient does not depend on the answer: the answer-dependent part
is $5\times10^{-3}$ ($n=2$) to $2\times10^{-5}$ ($n=10$) of the full gradient.}
\label{tab:spectrum}
\label{tab:spectrum-full}
\end{table}

\subsection{Prediction 4: outcome training plateaus}
\label{app:exp-plateau}
\label{app:basin}
\label{app:tracking}

\paragraph{Outcome training stays at chance across horizons.} \Cref{tab:gap} gives three seeds per
horizon. At $n=16$, with seeds $42$--$51$, step supervision ($K=2$) reaches $100\%$ exact continuation by
$16$k on nine of ten seeds, and outcome stays at chance on nine of ten at $32$k. At $K=3$, $n=16$,
outcome stays at chance on eight of ten; after $32$k, seeds $42$--$44$ give step
exact continuation $100\%$ each and outcome answer accuracy $50.0$, $33.5$ and $30.1\%$ (chance $33\%$). The $K=2$, $n=8$ runs in
\cref{tab:gap} have $49$--$55\%$ accuracy at $8$k and none of seeds $42$--$44$ reaches the
$90\%$ solved threshold by $32$k. A separate $K=3$, $n=8$ experiment, with
$\mixrate^8\approx0.043$, ends at $100.0\%$, $67.7\%$ and $99.9\%$ after $32$k for seeds
$42$, $43$ and $44$, respectively. Two of these modulo-$3$ runs solve the task.

\begin{table}[htbp]
\centering\small
\begin{tabular}{@{}lcccccc@{}}
\toprule
& & \multicolumn{2}{c}{step} & \multicolumn{2}{c}{outcome} &\\
\cmidrule(lr){3-4}\cmidrule(lr){5-6}
Task & $n$ & after $8$k & solved by $16$k & after $8$k & solved by $32$k & chance\\
\midrule
count, $K=2$ & $8$ & $57$--$95$ & $3/3$ & $49$--$55$ & $0/3$ & $50$\\
count, $K=2$ & $12$ & $100$ & $3/3$ & $49$--$51$ & $0/3$ & $50$\\
count, $K=2$ & $16$ & $17$--$26$ & $3/3$ & $49$--$52$ & $0/3$ & $50$\\
count, $K=2$ & $24$ & $74$--$100$ & $3/3$ & $49$--$52$ & $0/3$ & $50$\\
first occurrence & $8$--$24$ & $80$--$99$ & $12/12$ & $71$--$100$ & $11/12^\dagger$ & --\\
\bottomrule
\multicolumn{7}{@{}l}{\footnotesize $^\dagger$ solved by $16$k.}
\end{tabular}
\caption{\textbf{Modulo-$2$ outcome runs remain unsolved at the tested horizons.} Exact-continuation
accuracy (\%) on held-out words, seeds $42$--$44$. ``Solved'': runs reaching $\ge90\%$ by the stated
budget. Both handcoded executors are exact on held-out words at every horizon.}
\label{tab:gap}
\end{table}


\paragraph{Across $120$ runs, outcome solves none.} Counting with $K=2$, $8{,}000$ updates, the same
$\theta_0$, data and batches for both losses per seed, $20$ seeds per horizon. A run is solved if its
final exact continuation is at least $90\%$. ($n=4$ is excluded because it has fewer distinct
prompts than the training set.)
\begin{center}\small
\begin{tabular}{@{}lcccccc@{}}
\toprule
$n$ & $6$ & $8$ & $10$ & $12$ & $14$ & $16$\\
\midrule
step, solved of $20$ & $5$ & $8$ & $13$ & $13$ & $13$ & $8$\\
outcome, solved of $20$ & $0$ & $0$ & $0$ & $0$ & $0$ & $0$\\
outcome, median final exact & $50.2$ & $50.8$ & $49.8$ & $50.5$ & $50.1$ & $50.4$\\
\bottomrule
\end{tabular}
\end{center}

\paragraph{A larger initialization helps step training but not outcome training.} Seeds $42$--$44$,
$8{,}000$ updates, exact-continuation accuracy. Step training becomes reliable at $\initscale=0.1$.
The $\initscale^6$ credit scale would be $5^6\approx1.6\times10^4$ larger than at $0.02$.
This is a qualitative comparison; the theorem does not certify these scales.
Outcome training stays at chance at $n=16$ for all tested scales.
At $\initscale=0.5$, both losses are unstable at this learning rate.

\begin{center}\small
\begin{tabular}{@{}lccc@{}}
\toprule
& $\initscale=0.02$ & $\initscale=0.1$ & $\initscale=0.5$\\
\midrule
outcome, $n=8$ & $49$--$55\%$ & $63$--$65\%$ & $48$--$53\%$\\
outcome, $n=16$ & $48$--$52\%$ & $48$--$50\%$ & $50$--$52\%$\\
step, $n=8$ & $57$--$95\%$ & $100\%$ & $8$--$69\%$\\
step, $n=16$ & $17$--$26\%$ & $100\%$ & $0$--$1\%$\\
\bottomrule
\end{tabular}
\end{center}

\paragraph{Depth does not change the split.} With $\blocks\in\{4,6\}$ blocks, $n\in\{8,12\}$ and seeds
$42$--$44$, step training reaches $31$--$100\%$ after $8$k updates ($10$ of $12$ runs above $80\%$),
outcome training stays at $47.8$--$50.9\%$, and first occurrence reaches $95$--$100\%$ under both.

\paragraph{On the plateau the Jacobian stays low-degree.} We compute exact letter spectra of the
outcome model through training ($n=8$, $K=2$). The two bound columns are the two terms of
\cref{lem:gradient-split}. Above-degree-$1$ leakage is $0.5\%$ of the low-degree term at $1$k and $3$k and
$7\%$ at $6$k, and dominates only after escape. Averaging over $s_0$ cancels a further factor of about
$700$ in the measured credit, which the per-$(s_0,q)$ bound does not use. For seeds $43$ and $44$, $\Jhigh/\Jtotal$ stays
at most $7.6\times10^{-7}$ at every checkpoint, apart from one transient spike of $7.5\times10^{-3}$ for
seed $43$ at $6$k, where its answer-dependent gradient stays $7.8\times10^{-4}$; accuracy stays at
$48$--$52\%$.

\begin{table}[htbp]
\centering\small
\fitwidth{\begin{tabular}{@{}lcccccc@{}}
\toprule
outcome updates & degree-0 share & $\Jhigh/\Jtotal$ & $\sqrt{(A_0+A_1)(J_0+J_1)}$ & $\sqrt{(1-\frac1K)\Jhigh}$ & $\|\pd{\outcredit}{\theta}\|$ & accuracy\\
\midrule
init & $0.9964$ & $2.4\mathrm{e}{-6}$ & $5.0\mathrm{e}{-3}$ & $4.0\mathrm{e}{-4}$ & $2.1\mathrm{e}{-5}$ & --\\
$1{,}000$ & $0.9995$ & $1.1\mathrm{e}{-8}$ & $0.56$ & $3.0\mathrm{e}{-3}$ & $8.8\mathrm{e}{-5}$ & $52\%$\\
$3{,}000$ & $0.9995$ & $1.1\mathrm{e}{-8}$ & $0.41$ & $2.2\mathrm{e}{-3}$ & $1.6\mathrm{e}{-5}$ & $48\%$\\
$6{,}000$ & $0.9973$ & $2.0\mathrm{e}{-6}$ & $0.30$ & $2.2\mathrm{e}{-2}$ & $4.5\mathrm{e}{-4}$ & $48\%$\\
$8{,}000$ & $0.711$ & $0.155$ & $0.25$ & $5.5$ & $7.3\mathrm{e}{-2}$ & $55\%$\\
\bottomrule
\end{tabular}}
\caption{\textbf{On the plateau the outcome model has almost no energy above degree $1$.} Counting,
$K=2$, $n=8$, seed $42$.}
\label{tab:plateau-h1}
\end{table}


\subsection{Other tasks and pretrained models}
\label{app:exp-other}
\label{app:circuits}

GPT-2L uses character tokens.
Boolean-8 samples a gate type uniformly, then its wires.
Rank-8 LoRA on SmolLM2-135M and Qwen2.5-0.5B gives the same supervision gap.
Outcome training remains poor at eight times the budget.
On GPT-2L, answer-first and filler formats give $8.0\pm0.9\%$ and $7.8\pm0.5\%$
over five seeds. Step supervision gives $82.0\pm7.7\%$.
Extra tokens alone do not reproduce the benefit.

\begin{table}[htbp]
\centering\small
\begin{tabular}{@{}llccc@{}}
\toprule
Task & Condition & GPT-2L & SmolLM2-135M & Qwen2.5-0.5B\\
\midrule
Boolean-8 & outcome & $\msd{7.7}{0.3}$ & $\msd{9.9}{0.2}$ & $\msd{11.1}{0.2}$\\
 & outcome, $8\times$ budget & -- & $\msd{11.8}{2.2}$ & $\msd{10.0}{0.9}$\\
 & answer first & $\msd{8.0}{0.9}$ & $\msd{9.3}{0.3}$ & $\msd{10.1}{1.0}$\\
 & step & $\msd{82.0}{7.7}$ & $\msd{89.7}{4.9}$ & $\msd{98.3}{2.3}$\\
\midrule
register machine & outcome & $\msd{6.3}{0.8}$ & -- & --\\
 & step & $\msd{99.9}{0.2}$ & -- & --\\
\bottomrule
\end{tabular}
\caption{\textbf{Pretrained models repeat the split.} Answer accuracy (\%), mean$_{\pm\mathrm{SD}}$
over seeds. Chance is $6.25\%$ for Boolean-8 and $5.88\%$ for the register machine.}
\label{tab:lora}
\label{tab:further}
\end{table}

\subsection{BIG-Bench Hard object tracking}
\label{app:bbh}

\paragraph{Data.} Generate training items using the benchmark's templates, five scenes and item lists.
Sample new assignments and swaps, uniformly over pairs.
Test on the official $250$ items per size; none is used for training.
All $750$ items parse and rebuild exactly.

Zero-shot accuracy is $33.6$, $16.4$ and
$11.2\%$ for $3$, $5$ and $7$ objects, at or below chance, scored by the log-probability of each option
letter. Accuracy on real items tracks that on held-out generated items within about ten points.

\paragraph{Traces.} Each state line lists every holder's item and its option letter, e.g.\
``Alice has Ulysses (A)''. With item names alone, traces were correct but answers stayed at the letter base rate.
We therefore include option letters in the traces.

\paragraph{Low-degree share.} With $m$ objects, the item asked about has $K=m$ states, its holder.
A random swap keeps the holder with probability $p=\binom{m-1}{2}/\binom m2$ and moves it to each other
holder with probability $(1-p)/(m-1)$, so $\mixrate=(mp-1)/(m-1)$, which is $0$, $\frac12$ and $\frac23$ for $m=3,5,7$. The
share of answer variance at degrees $\le1$ after $n$ swaps is $\mixrate^{2n}+n\mixrate^{2(n-1)}(1-\mixrate^2)$.

\begin{table}[htbp]
\centering\small
\begin{tabular}{@{}lcccc@{}}
\toprule
objects (swaps) & degree-$\le1$ share & outcome, three runs & step, three runs & chance\\
\midrule
$3$ ($3$) & $0$ & $35.2$, $35.2$, $35.2$ & $100$, $98.4$, $100$ & $33$\\
$5$ ($5$) & $1.56\%$ & $94.4$, $83.6$, $80.0$ & $100$, $100$, $46.0$ & $20$\\
$7$ ($7$) & $3.34\%$ & $24.4$, $16.0$, $15.2$ & $96.4$, $100$, $100$ & $14$\\
\bottomrule
\end{tabular}
\caption{\textbf{BIG-Bench Hard object tracking on real items.}
Qwen2.5-0.5B, rank-8 LoRA. Runs share training data and order; only initialization changes.
At three objects, outcome runs select the most common option, ``(B)''.
Its base rate is $35.2\%$ on real items and $28.7\%$ on generated test items.
The failed five-object step run also scores $45\%$ on generated items.}
\label{tab:bbh}
\end{table}

\begin{table}[htbp]
\centering\small
\begin{tabular}{@{}lcccccc@{}}
\toprule
objects, swaps & $5$, $3$ & $7$, $6$ & $5$, $6$ & $7$, $10$ & $5$, $10$ & $3$, $10$\\
\midrule
degree-$\le1$ share & $15.6\%$ & $6.55\%$ & $0.46\%$ & $0.41\%$ & $0.003\%$ & $0$\\
outcome & $94.0$ & $77.0$ & $18.7$ & $15.7$ & $21.7$ & $39.7$--$40.7$\\
step & $100$ & $100$ & $100$ & $98.0$ & $100$ & $100$\\
chance & $20$ & $14$ & $20$ & $14$ & $20$ & $33$\\
\bottomrule
\end{tabular}
\caption{\textbf{At a fixed number of swaps the low-degree share orders outcome learning.} Pilot
on a plain template, same model and budget, accuracy (\%) on held-out generated items, one seed
(three for three objects). Outcome learns where the share is at least $6.5\%$ and stays near chance,
or the majority rate ($35.7\%$ at three objects), where it is at most $0.5\%$.}
\label{tab:bbh-pilot}
\end{table}

%=============================================================================
\section{Reproducibility}
\label{app:reproducibility}

Run \texttt{python3} on \path{handcoded/lookup_convergence.py} for the fixed-feature check.
It uses NumPy in float64 and saves losses, barrier-crossing updates and bound checks under
\path{results/lookup_convergence/check.json}.

Code, configurations and seeds are at \url{https://github.com/anonymous/trace}.
\texttt{RUN.md} lists the experiment commands.
Training data are synthetic.
BIG-Bench Hard test items come from the official repository
\citep{suzgun2022challenging} and are stored unchanged in \texttt{data/bbh}.
Other test sets are generated.

Software and models, with licenses: PyTorch $2.14$
\citep{paszke2019pytorch} (BSD-3-Clause), Hugging Face Transformers $5.16$ and PEFT $0.20$
(Apache-2.0), SmolLM2-135M and Qwen2.5-0.5B (Apache-2.0), and the BIG-Bench Hard items (MIT).
Compute: a shared node with four NVIDIA A100 80GB GPUs, one GPU per run. A letter-task run of
$8{,}000$ updates takes about one minute, and a Qwen2.5-0.5B LoRA run of $800$ updates with
evaluation about fifteen minutes. Exploratory runs used more compute than the reported
experiments.
